\documentclass[11pt,a4paper]{report}

% ===================== Encodage & langue =====================
\usepackage[T1]{fontenc}
\usepackage[utf8]{inputenc}
\usepackage{lmodern}
\usepackage[french]{babel}
\usepackage{amsmath}
\usepackage{amssymb}

% ===================== Mise en page =====================
\usepackage[a4paper,margin=2.5cm,headheight=14pt]{geometry}
\usepackage{fancyhdr}
\usepackage{setspace}
\onehalfspacing
\usepackage{parskip}
\usepackage{microtype}

% ===================== Graphiques & tableaux =====================
\usepackage{graphicx}
\usepackage{xcolor}
\usepackage{booktabs}
\usepackage{array}
\usepackage{longtable}
\usepackage{multirow}
\usepackage{enumitem}
\usepackage{tikz}
\usetikzlibrary{shapes.geometric,shapes.multipart,arrows.meta,positioning,calc,fit,backgrounds,trees}
\usepackage{pgfplots}
\pgfplotsset{compat=1.18}
\usepackage{pgfgantt}

% ===================== Code source (listings) =====================
\usepackage{listings}

\definecolor{codebg}{rgb}{0.97,0.97,0.95}
\definecolor{codekw}{rgb}{0.10,0.10,0.65}
\definecolor{codecmt}{rgb}{0.30,0.55,0.30}
\definecolor{codestr}{rgb}{0.65,0.15,0.15}
\definecolor{codeframe}{rgb}{0.80,0.80,0.80}
\definecolor{linenoclr}{rgb}{0.55,0.55,0.55}

% Rendu robuste de l'UTF-8 dans les listings (accents, guillemets...).
\lstset{
  basicstyle=\ttfamily\footnotesize,
  backgroundcolor=\color{codebg},
  keywordstyle=\color{codekw}\bfseries,
  commentstyle=\color{codecmt}\itshape,
  stringstyle=\color{codestr},
  numberstyle=\tiny\color{linenoclr},
  numbers=left,
  numbersep=8pt,
  frame=single,
  rulecolor=\color{codeframe},
  breaklines=true,
  breakatwhitespace=false,
  showstringspaces=false,
  tabsize=4,
  captionpos=b,
  extendedchars=true,
  literate=%
    {é}{{\'e}}1 {è}{{\`e}}1 {ê}{{\^e}}1 {ë}{{\"e}}1
    {à}{{\`a}}1 {â}{{\^a}}1 {ä}{{\"a}}1
    {ù}{{\`u}}1 {û}{{\^u}}1 {ü}{{\"u}}1
    {î}{{\^i}}1 {ï}{{\"i}}1
    {ô}{{\^o}}1 {ö}{{\"o}}1
    {ç}{{\c c}}1 {Ç}{{\c C}}1
    {É}{{\'E}}1 {È}{{\`E}}1 {Ê}{{\^E}}1
    {À}{{\`A}}1
    {«}{{\guillemotleft{}}}1 {»}{{\guillemotright{}}}1
    {’}{{'}}1 {…}{{...}}1 {–}{{-}}1 {—}{{---}}1
    {°}{{\textdegree}}1
}

% Style pour le C.
\lstdefinestyle{c}{language=C}
% Style "grammaire" (Flex/Bison) : pas de langage predefini, on garde la
% coloration des chaines et commentaires.
\lstdefinestyle{gram}{
  morecomment=[l]{//},
  morecomment=[s]{/*}{*/},
  morestring=[b]",
}
% Style assembleur AT&T.
\lstdefinestyle{asm}{
  morecomment=[l]{\#},
  morestring=[b]",
}
% Style ILANG.
\lstdefinestyle{ilang}{
  morekeywords={si,sinon,tant,que,pour,fonction,retourner,afficher,lire,arreter,continuer},
  morecomment=[l]{//},
  morecomment=[s]{/*}{*/},
}

% ===================== Liens (charge en dernier) =====================
\usepackage[hidelinks,bookmarks=true,bookmarksnumbered=true]{hyperref}

% ===================== En-tetes / pieds de page =====================
\pagestyle{fancy}
\fancyhf{}
\fancyhead[L]{\small AIT BENAISSA Ilias --- 82508880}
\fancyhead[R]{\small Interpréteur ILANG --- IC (Chap. 10)}
\fancyfoot[C]{\thepage}
\renewcommand{\headrulewidth}{0.4pt}
\renewcommand{\footrulewidth}{0pt}

% Raccourcis.
\newcommand{\ilang}{\textsc{Ilang}}
\newcommand{\code}[1]{\texttt{#1}}

\begin{document}

% ===================== Page de titre =====================
\begin{titlepage}
\centering
\vspace*{1.5cm}
{\scshape\large Institut d'Enseignement à Distance --- Université Paris 8\par}
{\scshape Licence 3 Informatique\par}
\vspace{0.5cm}
{\large Interprétation \& Compilation --- Chapitre 10 : Projet Final\par}
\vspace{2.5cm}
\rule{\linewidth}{0.5mm}\par
\vspace{0.8cm}
{\Huge\bfseries Conception et implémentation\\[0.3cm] de l'interpréteur \ilang\par}
\vspace{0.6cm}
{\Large\itshape Un mini-langage impératif réalisé avec Flex et Bison,\\
de l'analyse lexicale à la génération de code\par}
\vspace{0.8cm}
\rule{\linewidth}{0.5mm}\par
\vspace{2.5cm}
{\large\bfseries AIT BENAISSA Ilias\par}
{\large N\textsuperscript{o} étudiant : 82508880\par}
\vspace{1.5cm}
{\large Enseignant : Philippe Kislin-Duval\par}
\vfill
{\large \today\par}
\end{titlepage}

% ===================== Resume =====================
\chapter*{Résumé}
\addcontentsline{toc}{chapter}{Résumé}
Ce rapport présente la conception et l'implémentation complète d'un interpréteur
pour \ilang, un mini-langage de programmation impératif conçu dans le cadre du
projet final du cours d'Interprétation \& Compilation. Le langage supporte les
variables, les expressions arithmétiques avec priorité des opérateurs, les
structures de contrôle (conditions, boucles), les fonctions utilisateur
récursives et les entrées-sorties. Sa réalisation réinvestit l'ensemble de la
chaîne de compilation étudiée en cours : analyse lexicale par automate fini
(Flex), analyse syntaxique par grammaire LALR(1) (Bison), construction d'un arbre
syntaxique abstrait (AST), optimisations sur cet arbre (pliage de constantes,
élimination du code mort), puis évaluation récursive.

Au-delà du cahier des charges, le projet a été prolongé par trois extensions
(nombres flottants, instructions \code{arrêter}/\code{continuer}, localisation
des erreurs à la ligne) et par une exploration hors-périmètre particulièrement
formatrice : un \emph{second back-end} qui, à partir du même AST, génère du code
assembleur i386 32~bits compilable avec \code{gcc -m32}. Cette dualité
interpréteur/compilateur illustre concrètement le principe d'un front-end commun
suivi de back-ends interchangeables.

L'implémentation, écrite en C, compile sans aucun avertissement ni conflit
grammatical et ne présente aucune fuite mémoire (vérifié avec Valgrind). Une
campagne de tests systématique, couvrant les programmes de démonstration, la
gestion des erreurs, les cas limites et les mesures de performance, valide le bon
fonctionnement et
met en évidence les forces et les limites de l'approche, notamment la profondeur
de récursion bornée par la pile de l'hôte et l'écart de performance considérable
entre l'interprétation et le code natif généré.

\vspace{0.5cm}
\noindent\textbf{Mots-clés :} interpréteur, compilation, Flex, Bison, LALR(1),
arbre syntaxique abstrait, table des symboles, pliage de constantes, génération
de code, assembleur i386.


% ===================== Table des matieres =====================
\tableofcontents
\listoffigures
\lstlistoflistings

% ===================== Chapitres =====================
\chapter{Introduction et contexte}
\label{chap:intro}

\section{Contexte et problématique}

\subsection{Contexte pédagogique et scientifique}

Ce projet constitue l'évaluation finale du cours d'Interprétation et Compilation
de la Licence~3 Informatique de l'IED, à l'Université Paris~8. Il occupe une
place particulière dans la progression du semestre, car il ne porte pas sur une
notion isolée mais demande au contraire de rassembler en un seul logiciel
fonctionnel l'ensemble des concepts étudiés tout au long de l'année. Les chapitres
précédents avaient abordé séparément l'assembleur et la pile d'appels, la chaîne
de compilation, les automates finis, les expressions régulières, les grammaires,
puis les optimisations. Le projet final consiste à montrer qu'on est capable de
faire tenir tous ces éléments ensemble, de les articuler les uns aux autres, et
d'obtenir au bout un programme qui exécute réellement du code écrit dans un langage
que l'on a soi-même défini.

Il faut rappeler à quel point la traduction des langages occupe une position
centrale en informatique. Aucun programme écrit dans un langage de haut niveau
n'est directement compréhensible par le processeur, qui ne sait exécuter que des
instructions machine élémentaires. Il existe donc toujours, quelque part, un
logiciel chargé de faire le pont entre ce que le programmeur écrit et ce que la
machine exécute. Deux grandes familles de stratégies remplissent ce rôle. La
compilation traduit le programme une fois pour toutes en code machine, qui sera
ensuite exécuté de façon autonome. L'interprétation, à l'inverse, exécute le
programme directement, en parcourant une représentation interne, sans produire de
binaire séparé. Ces deux familles, malgré leur différence de finalité, partagent
exactement la même première moitié de traitement : analyser le texte source, le
décomposer en unités lexicales, vérifier sa conformité à une grammaire, et le
représenter sous une forme structurée. Elles ne divergent qu'ensuite, sur la
manière d'exploiter cette représentation. C'est précisément cette architecture en
deux temps, un tronc commun suivi de traitements spécialisés, que mon projet met
en évidence, d'autant plus nettement qu'il en propose deux exploitations
différentes du même tronc.

\subsection{Identification du problème}

Le problème posé peut s'énoncer simplement : concevoir un langage de programmation
complet, si modeste soit-il, et écrire le logiciel capable de l'exécuter. Le mot
« complet » mérite une précision. Il ne signifie pas riche en fonctionnalités,
mais suffisant pour exprimer n'importe quel calcul, c'est-à-dire Turing-complet.
En pratique, cela impose au minimum de disposer de variables pour mémoriser des
valeurs, de conditionnelles pour choisir entre plusieurs comportements, de boucles
pour répéter des traitements, et d'un mécanisme de composition, ici les fonctions,
pour structurer et réutiliser le code.

Concevoir un tel langage ne se réduit pas à écrire du code C. C'est un travail
préalable de définition qui porte sur trois aspects indissociables, et c'est
seulement une fois ces trois aspects fixés que l'écriture de l'interpréteur prend
son sens. Le premier aspect est la syntaxe concrète, c'est-à-dire l'apparence
qu'aura un programme valide : quels mots-clés, quels opérateurs, quelle
ponctuation. Le deuxième aspect est la grammaire, c'est-à-dire la définition
précise des suites de symboles acceptées et de la manière dont elles s'organisent
en structures imbriquées. Le troisième aspect, enfin, est la sémantique,
c'est-à-dire la signification de chaque construction, ce que fait réellement le
programme lorsqu'on l'exécute. La syntaxe relève de Flex, la grammaire de Bison,
et la sémantique de l'évaluateur que j'ai écrit.

\subsection{Questions directrices}

Plutôt que d'aborder le projet comme une simple commande à satisfaire, j'ai
préféré l'organiser autour de quelques questions concrètes d'ingénierie, qui ont
guidé mes choix et qui structurent ce rapport. La première question concerne les
outils : qu'apporte réellement le recours à Flex et Bison par rapport à une analyse
écrite à la main, et ce gain justifie-t-il la courbe d'apprentissage de ces
outils ? La deuxième question concerne la représentation interne : comment
représenter un programme en mémoire de manière à pouvoir aussi bien l'évaluer,
l'optimiser, voire le compiler, sans avoir à tout réécrire pour chaque usage ? La
troisième question touche à la sémantique des fonctions : comment gérer
correctement la portée des variables et la récursion, qui sont les points les plus
délicats d'un langage à fonctions ? La dernière question, plus ouverte, est celle
des limites : jusqu'où peut-on pousser un projet de ce type avant de buter sur des
obstacles de fond, qu'il s'agisse de performance, de profondeur de récursion ou de
typage ?

\section{Objectifs et contributions}

\subsection{Objectifs}

L'objectif principal est de livrer un interpréteur \ilang{} qui soit à la fois
fonctionnel, robuste et fidèle au cahier des charges. Cet objectif se décompose en
plusieurs sous-objectifs concrets. Il faut d'abord écrire un analyseur lexical qui
reconnaisse toutes les unités lexicales du langage, des nombres aux mots-clés en
passant par les commentaires. Il faut ensuite écrire une grammaire LALR(1) qui
soit non ambiguë et qui construise au fil de l'analyse l'arbre syntaxique. Il faut
implémenter un évaluateur capable de traiter les variables, les expressions, les
structures de contrôle, les fonctions récursives et les entrées-sorties. Il faut
mettre en place une table des symboles correcte du point de vue de la portée. Il
faut réaliser les deux passes d'optimisation demandées. Il faut traiter proprement
les erreurs courantes en les localisant à la ligne. Enfin, et ce dernier point a
compté autant que les autres à mes yeux, il faut atteindre un certain niveau de
qualité d'ingénierie : une compilation sans le moindre avertissement, une grammaire
sans conflit, et une exécution sans fuite mémoire.

\subsection{Contributions}

Au-delà de ce que demandait strictement la commande, mon travail apporte trois
contributions que je revendique. La première est un ensemble de trois extensions
du langage. J'ai ajouté un type flottant, conçu de façon à cohabiter avec les
entiers sans compromettre la sémantique entière d'origine, ce qui était la vraie
difficulté. J'ai ajouté les instructions de rupture de boucle \code{arrêter} et
\code{continuer}. J'ai enfin systématisé la localisation des erreurs à la ligne.
La deuxième contribution est un second back-end, qui génère du code assembleur
i386 à partir du même arbre syntaxique. Je l'ai développé par curiosité, en sachant
qu'il dépassait le périmètre du projet, mais il démontre concrètement que le
front-end que j'avais écrit était réutilisable tel quel pour produire du code
natif, ce qui valide a posteriori toute l'architecture. La troisième contribution
est une démarche de validation que j'ai voulue sérieuse, comprenant des tests
fonctionnels, des tests d'erreurs, une campagne de cas limites, une analyse de la
mémoire et des mesures de performance comparant l'interprétation au code natif.

\subsection{Cadrage du projet par les cinq questions}

Pour cadrer clairement le travail, il est utile de répondre explicitement aux cinq
questions classiques de la conduite de projet, souvent appelées les cinq W. La
question du \emph{quoi} a déjà été posée : il s'agit d'un interpréteur pour le
mini-langage impératif \ilang, accompagné de ses programmes de démonstration et de
sa documentation. La question du \emph{pourquoi} renvoie à la finalité
pédagogique : il s'agit de démontrer la maîtrise de la chaîne de compilation
complète, ce qui est l'objectif assigné au projet final du chapitre~10. La question
du \emph{qui} est vite réglée, puisqu'il s'agit d'un travail individuel évalué,
réalisé par un seul étudiant. La question du \emph{quand} renvoie à la période
allouée en fin de semestre, dont la répartition est présentée sous forme de
diagramme de Gantt au chapitre~\ref{chap:conception}. La question du \emph{où},
enfin, désigne l'environnement de travail, à savoir un système GNU/Linux exécuté
sous WSL, avec la chaîne d'outils GCC, Flex et Bison.

\subsection{Budget, temps et coûts}

S'agissant d'un projet académique individuel, la question financière se règle en
une phrase : le coût direct est nul. Tous les outils que j'ai employés, du
compilateur GCC aux générateurs Flex et Bison, en passant par Make, Valgrind et la
distribution \TeX{} Live qui a servi à produire ce rapport, sont des logiciels
libres et gratuits. La seule ressource réellement consommée est le temps. J'ai
réparti cet effort de manière progressive, en commençant par la conception de la
grammaire, puis l'écriture du c\oe ur du langage, ensuite la table des symboles et
les optimisations, puis une phase substantielle de débogage et de tests, et enfin
les extensions avant la rédaction. Cette répartition est détaillée et justifiée au
chapitre~\ref{chap:conception}, où elle est présentée sous forme de planning.

\section{Réinvestissement des acquis du cours}

Le projet est explicitement conçu comme une synthèse, et il me semble important de
rendre cette intention visible dès l'introduction. Le tableau~\ref{tab:reinvest}
met en regard chaque chapitre du cours et la manière précise dont il se retrouve
mobilisé dans \ilang. Cette mise en correspondance n'est pas un simple exercice de
style. Elle constitue un fil conducteur que je m'efforce de suivre tout au long du
rapport, en signalant, à chaque fois que c'est pertinent, à quelle notion du cours
se rattache tel choix d'implémentation.

\begin{table}[ht]
\centering
\caption{Réinvestissement des chapitres du cours dans le projet.}
\label{tab:reinvest}
\small
\begin{tabular}{@{}p{4.8cm}p{9cm}@{}}
\toprule
\textbf{Chapitre du cours} & \textbf{Concept réinvesti dans \ilang} \\
\midrule
2, Assembleur & Pile d'appels et cadre \code{\%ebp}, reproduits pour les fonctions,
et générés tels quels par le back-end assembleur. \\
3, Programme C compilé & Niveaux d'optimisation et pliage de constantes. \\
4, Analyse lexicale et syntaxique & Analyseur lexical écrit avec Flex et parseur
écrit avec Bison. \\
5, Yacc/Bison & Grammaire LALR(1) et résolution, ou plutôt constat d'absence, des
conflits shift/reduce. \\
7, Génération de code & Arbre syntaxique abstrait, table des symboles, évaluation
récursive. \\
8, Optimisations & Pliage de constantes et élimination du code mort sur l'arbre. \\
9, Automates et regex & Tokens reconnus par l'automate fini engendré par Flex. \\
\bottomrule
\end{tabular}
\end{table}

\section{Structure du document}

Le rapport suit une progression que je veux la plus naturelle possible, allant du
général au particulier puis à la prise de recul. Le chapitre~\ref{chap:etat} dresse
un état de l'art des techniques d'interprétation et de compilation et explique
comment mon travail se positionne par rapport à elles. Le
chapitre~\ref{chap:conception} présente l'analyse des besoins, l'architecture que
j'ai retenue et la conception détaillée, en particulier la grammaire, l'arbre
syntaxique et les algorithmes principaux. Le chapitre~\ref{chap:implementation}
décrit l'implémentation proprement dite, module par module, avec de nombreux
extraits de code commentés et des exemples concrets. Le
chapitre~\ref{chap:validation} expose ma stratégie de test et les résultats
obtenus, y compris les mesures de performance et une étude de cas de bout en bout.
Le chapitre~\ref{chap:discussion} prend du recul, analyse de façon critique les
forces et les limites du projet, et revient en détail sur les difficultés que j'ai
rencontrées. Le chapitre~\ref{chap:conclusion} conclut et ouvre des perspectives.
Les annexes regroupent enfin l'intégralité du code source, les manuels
d'installation et d'utilisation, l'assembleur réellement généré et les données
expérimentales.

\chapter{État de l'art}
\label{chap:etat}

Concevoir un langage et l'outil capable de l'exécuter n'est pas une démarche
isolée. C'est au contraire s'inscrire dans une tradition technique et théorique
qui s'est constituée sur plus d'un demi-siècle, et qui a fini par dégager des
méthodes et des outils devenus standards. Avant de présenter mes propres choix,
il me paraît donc indispensable de situer le projet dans ce paysage : de rappeler
d'où viennent les concepts que j'ai utilisés, quelles solutions existent pour
chaque problème que j'ai eu à résoudre, et pourquoi j'ai retenu telle approche
plutôt que telle autre. C'est l'objet de ce chapitre, organisé en une revue des
travaux et des outils existants, suivie du positionnement de mon travail.

\section{Revue de littérature}

\subsection{Travaux fondateurs}

Tout ce projet repose, souvent sans qu'on y pense, sur des résultats théoriques
établis entre les années 1950 et 1970. Le premier de ces apports est la
classification des grammaires formelles proposée par Noam Chomsky en
1956~\cite{chomsky}. Chomsky a montré qu'on peut hiérarchiser les langages selon
la puissance de la machine nécessaire pour les reconnaître. Deux niveaux de cette
hiérarchie nous intéressent directement. Les langages dits réguliers, les plus
simples, sont reconnus par des automates finis : ils suffisent à décrire la forme
des mots du langage, c'est-à-dire les nombres, les identifiants ou les
opérateurs. Les langages dits hors-contexte, plus expressifs, sont reconnus par
des automates à pile : ils sont nécessaires pour décrire l'imbrication des
parenthèses, des blocs et des expressions. Cette distinction n'est pas une
subtilité académique. Elle se traduit très concrètement dans l'architecture de
n'importe quel compilateur, et dans la mienne en particulier, par la séparation
en deux étages bien distincts. L'analyse lexicale, confiée à Flex, traite le
niveau régulier. L'analyse syntaxique, confiée à Bison, traite le niveau
hors-contexte. J'ai trouvé éclairant de constater que cette frontière théorique
correspond exactement à la frontière entre mes deux premiers fichiers source.

Le deuxième apport fondateur concerne l'analyse syntaxique elle-même. En 1965,
Donald Knuth a décrit l'analyse LR~\cite{knuth}, une méthode ascendante,
déterministe et efficace, capable de reconnaître une très large classe de
grammaires hors-contexte. Le mot « ascendante » signifie que l'analyseur part des
symboles élémentaires et reconstruit progressivement les structures
englobantes, par opposition aux méthodes descendantes qui partent du but et le
décomposent. Les variantes plus économes en mémoire de cette méthode, en
particulier l'analyse LALR(1) qui n'utilise qu'un seul symbole de prévision, sont
précisément celles que met en \oe uvre l'outil Bison. Autrement dit, lorsque
j'écris une grammaire pour Bison, je m'appuie sans le réécrire sur un algorithme
qui a près de soixante ans et qui fait toujours référence.

L'ensemble de cette chaîne, de l'analyse lexicale jusqu'à la génération de code en
passant par les arbres syntaxiques et les optimisations, a été synthétisé et
codifié dans un ouvrage devenu la référence absolue du domaine, communément appelé
le « Dragon Book », d'Aho, Lam, Sethi et Ullman~\cite{dragon}. La structure même
de mon projet, et donc de ce rapport, suit fidèlement le découpage proposé par cet
ouvrage. Je m'y suis appuyé en particulier pour comprendre le rôle pivot de
l'arbre syntaxique abstrait et la notion de table des symboles.

\subsection{Des outils générateurs plutôt que du code écrit à la main}

Sur le plan des outils, l'apport décisif date de 1975, aux laboratoires Bell, avec
deux programmes qui ont changé la manière d'écrire les compilateurs. Le premier,
\texttt{lex}, dû à Mike Lesk et Eric Schmidt~\cite{lesk}, engendre un analyseur
lexical à partir d'une simple liste d'expressions régulières. Le second,
\texttt{yacc} pour \emph{Yet Another Compiler-Compiler}, dû à Stephen
Johnson~\cite{johnson}, engendre un analyseur syntaxique à partir d'une grammaire.
L'idée commune à ces deux outils est qu'on ne devrait pas écrire un analyseur à la
main, instruction par instruction, mais le \emph{décrire} de façon déclarative et
laisser une machine produire le code correspondant. Cette idée s'est imposée
durablement, et les réimplémentations libres de ces outils, Flex~\cite{flex} pour
l'analyse lexicale et Bison~\cite{bison} pour l'analyse syntaxique, sont
aujourd'hui omniprésentes. Ce sont elles que j'ai employées. L'ouvrage de John
Levine~\cite{levine} m'a servi de guide pratique pour leur utilisation conjointe,
notamment pour comprendre comment le lexer et le parseur communiquent à travers la
variable \texttt{yylval} et la table des tokens.

\subsection{Les stratégies d'exécution}

Une fois le programme analysé et représenté en mémoire, encore faut-il l'exécuter.
Plusieurs stratégies existent, et le choix de l'une d'elles est probablement la
décision d'architecture la plus structurante d'un tel projet. Je les présente ici
de la plus simple à la plus performante, car cette progression correspond aussi à
une montée en complexité de réalisation.

La première stratégie est l'interprétation par parcours d'arbre, souvent désignée
par son nom anglais de \emph{tree-walking}. Le principe consiste à parcourir
récursivement l'arbre syntaxique et à évaluer chaque n\oe ud au moment où on le
rencontre. Un n\oe ud d'addition évalue ses deux fils puis additionne les
résultats, un n\oe ud de condition évalue son test puis descend dans la bonne
branche, et ainsi de suite. C'est l'approche la plus directe et, surtout, la plus
lisible : la sémantique du langage se lit presque mot pour mot dans le code de
l'évaluateur. Son inconvénient est la lenteur, car chaque exécution d'une
instruction suppose de re-parcourir une portion de l'arbre et de déterminer à
chaque n\oe ud de quel type il s'agit.

La deuxième stratégie repose sur une machine virtuelle à bytecode. Le programme
est d'abord traduit dans un jeu d'instructions compact et de bas niveau, le
bytecode, qui est ensuite exécuté par une boucle d'interprétation dédiée. C'est
l'approche retenue par des langages très répandus comme Python, Java ou Lua. Elle
est sensiblement plus rapide que le parcours d'arbre, car le travail d'analyse des
structures est fait une fois pour toutes au moment de la traduction, tout en
restant portable d'une machine à l'autre.

La troisième stratégie, la compilation à la volée ou \emph{JIT}, va plus loin en
traduisant en code machine, pendant l'exécution, les portions de programme les
plus sollicitées. C'est l'approche des moteurs JavaScript modernes ou de la
machine virtuelle Java dans sa version optimisée. Elle offre d'excellentes
performances au prix d'une complexité d'implémentation considérable. La quatrième
et dernière stratégie est la compilation anticipée vers du code machine, telle que
la pratiquent GCC ou Clang pour le langage C. Elle produit les binaires les plus
rapides, mais elle suppose une étape de compilation séparée et un résultat lié à
une architecture précise.

\subsection{Analyse comparative}

Le tableau~\ref{tab:approches} résume les compromis que présentent ces quatre
stratégies. Il ne faut pas le lire comme un classement, car aucune approche n'est
supérieure dans l'absolu : tout dépend de l'objectif poursuivi. Pour un produit
industriel soucieux de performance, le bytecode ou la compilation s'imposent. Pour
un projet dont la finalité est de \emph{comprendre} et de \emph{montrer} comment
un langage fonctionne, c'est au contraire la lisibilité qui prime, et le parcours
d'arbre devient le meilleur choix.

\begin{table}[ht]
\centering
\caption{Comparaison des stratégies d'exécution.}
\label{tab:approches}
\begin{tabular}{@{}lccc@{}}
\toprule
\textbf{Approche} & \textbf{Performance} & \textbf{Complexité} & \textbf{Lisibilité} \\
\midrule
Parcours d'arbre        & faible      & faible   & excellente \\
Machine à bytecode      & moyenne     & moyenne  & moyenne \\
Compilation JIT         & élevée      & élevée   & faible \\
Compilation anticipée   & très élevée & élevée   & faible \\
\bottomrule
\end{tabular}
\end{table}

Un point mérite d'être souligné, car il a orienté toute la suite du projet.
\ilang{} illustre en réalité deux points opposés de cet axe. La version
principale, l'interpréteur, relève du parcours d'arbre. L'extension que j'ai
développée par curiosité, le générateur de code assembleur, relève de la
compilation anticipée. Or les deux partent du \emph{même} arbre syntaxique. Cette
coïncidence n'en est pas une : elle matérialise le fait, central dans toute la
théorie de la compilation, que l'analyse d'un programme est indépendante de ce
qu'on décide d'en faire ensuite. C'est cette idée que j'ai trouvée la plus
instructive de tout le projet, et j'y reviendrai longuement.

\subsection{Les familles d'outils d'analyse syntaxique}

Puisque le c\oe ur de mon travail repose sur un générateur d'analyseur syntaxique,
il est utile de rappeler les grandes familles d'outils disponibles, afin de mieux
justifier ensuite le choix de Bison. La première famille regroupe les générateurs
ascendants de type LR ou LALR, dont Yacc et Bison sont les représentants les plus
connus. Ils construisent automatiquement une table d'analyse à partir de la
grammaire, reconnaissent une large classe de langages, et possèdent une vertu
précieuse pour qui apprend : ils signalent les ambiguïtés de la grammaire sous
forme de conflits. La deuxième famille regroupe les générateurs descendants de
type LL, comme ANTLR ou JavaCC. Ils sont souvent jugés plus intuitifs à déboguer,
car le déroulement de l'analyse suit visiblement la structure de la grammaire,
mais ils sont restreints à une classe de grammaires plus étroite et ont
longtemps été gênés par la récursion à gauche, qu'il faut éliminer à la main.

La troisième famille n'est pas un outil mais une technique : la descente récursive
écrite directement à la main, sans générateur. Elle est souple, lisible pour de
petits langages, et ne dépend d'aucun outil externe. Son défaut majeur, du point
de vue de ce projet, est que la grammaire y reste implicite, dispersée dans le
code, et que rien ne signale ses éventuelles ambiguïtés. La quatrième famille,
plus récente, regroupe les grammaires d'expression d'analyse, dites PEG, et les
combinateurs d'analyseurs. Elles fusionnent souvent l'analyse lexicale et
syntaxique et reposent sur un choix ordonné des règles, ce qui présente l'avantage
de ne jamais produire d'ambiguïté, mais au prix de la masquer plutôt que de la
révéler. Cette dernière propriété est exactement à l'opposé de ce que je
recherchais : je voulais un outil qui me \emph{dise} si ma grammaire était
ambiguë, pas un outil qui décide silencieusement à ma place.

\section{Positionnement}

\subsection{Lacunes de l'approche antérieure}

Le cahier des charges fait explicitement référence à un exercice antérieur du
cours, le mini-langage \textsc{Parlot}, qui réalisait son analyse lexicale au
moyen d'expressions régulières écrites en Python et son analyse syntaxique par une
descente récursive codée à la main. Cette approche fonctionne, et elle a le mérite
de la simplicité initiale, mais elle souffre de deux faiblesses que j'ai cherché à
corriger. La première est que la grammaire n'y est formalisée nulle part. Elle est
diluée dans le code impératif de la descente récursive, si bien qu'on ne peut ni
la lire d'un seul tenant, ni raisonner sur ses propriétés, ni détecter ses
ambiguïtés. La seconde faiblesse découle de la première : comme aucune description
formelle n'existe, rien ne vérifie que tous les cas sont traités. C'est au
programmeur seul d'y veiller, et la moindre ambiguïté passe inaperçue jusqu'à ce
qu'elle provoque un comportement inattendu.

\subsection{Opportunités et justification de mon approche}

Mon projet corrige ces deux lacunes en adoptant les outils standards du cours. En
décrivant la grammaire de façon déclarative dans le fichier Bison, j'obtiens deux
bénéfices que je considère comme décisifs. Le premier est que la grammaire devient
un objet explicite, écrit en notation BNF, que l'on peut lire, discuter, vérifier
et faire évoluer indépendamment du reste du code. Le second, plus important
encore, est que Bison détecte automatiquement les ambiguïtés sous forme de
conflits. Cette détection n'est pas restée théorique dans mon cas : elle m'a
permis de vérifier formellement que ma grammaire était non ambiguë, et même de
faire une découverte qui m'a marqué, l'absence du fameux conflit du « dangling
else », que je détaille au chapitre~\ref{chap:discussion}.

Reste à justifier le choix de l'interprétation par parcours d'arbre plutôt qu'une
machine à bytecode. Ce choix obéit au même principe directeur que le précédent, à
savoir privilégier la clarté. Une machine à bytecode aurait été plus rapide, mais
elle aurait introduit un niveau d'indirection supplémentaire, le jeu
d'instructions intermédiaire, qu'il aurait fallu concevoir, documenter et déboguer
pour un gain de performance sans intérêt dans un cadre pédagogique. Le parcours
d'arbre, au contraire, maintient une correspondance immédiate et visible entre la
structure du programme, représentée par l'arbre, et sa signification, codée dans
l'évaluateur. Ce choix s'est révélé pleinement justifié a posteriori : sa
régularité m'a permis d'ajouter trois extensions au langage, puis un générateur de
code complet, sans jamais avoir à remettre en cause l'architecture de départ.

\chapter{Analyse et conception}
\label{chap:conception}

\section{Analyse des besoins}

\subsection{Besoins fonctionnels}

Le langage \ilang{} et son interpréteur doivent offrir les fonctionnalités
suivantes, directement issues du cahier des charges.

\begin{description}[style=nextline,leftmargin=1em]
  \item[Types de données.] Entiers (seul type initialement requis) ; les booléens
        sont encodés en entiers (0 = faux, tout le reste = vrai), à la manière du
        C.
  \item[Variables.] Déclaration implicite à la première affectation, sans
        mot-clé. Portée globale pour le programme principal, locale pour les
        fonctions.
  \item[Expressions arithmétiques.] Opérateurs \code{+ - * / \%}, moins unaire,
        parenthèses, avec priorité et associativité correctes.
  \item[Comparaisons.] \code{== != < > <= >=}, renvoyant 0 ou 1.
  \item[Structures de contrôle.] Conditionnelle \code{si/sinon}, boucles
        \code{tant que} et \code{pour}.
  \item[Fonctions.] Déclaration avec \code{fonction}, paramètres, valeur de
        retour via \code{retourner}, appels récursifs.
  \item[Entrées-sorties.] \code{afficher(expr)} et \code{lire()}.
  \item[Commentaires.] De ligne (\code{//}) et de bloc (\code{/* */}).
\end{description}

\subsection{Besoins non-fonctionnels}

\begin{description}[style=nextline,leftmargin=1em]
  \item[Correction.] L'interpréteur doit produire les résultats attendus sur les
        programmes de démonstration et gérer proprement les erreurs.
  \item[Robustesse.] Aucune entrée, même fautive, ne doit provoquer de
        comportement indéfini silencieux ; les erreurs sont signalées et
        localisées.
  \item[Qualité du code.] Compilation sans avertissement (\code{-Wall -Wextra}),
        sans conflit grammatical, sans fuite mémoire.
  \item[Portabilité.] Code C standard, compilable sur toute plateforme POSIX
        disposant de GCC, Flex et Bison.
  \item[Lisibilité et maintenabilité.] Découpage en modules à responsabilité
        unique, code commenté.
\end{description}

\subsection{Cas d'utilisation}

L'acteur unique est l'utilisateur (ici, le développeur de programmes \ilang). Les
cas d'utilisation principaux sont : (1)~exécuter un fichier source
(\code{./ilang prog.ilang}) ; (2)~saisir un programme au clavier en mode
interactif (\code{./ilang}) ; (3)~en extension, produire l'assembleur
correspondant (\code{./ilang -s prog.ilang}). Dans tous les cas, le programme lit
le code source, le traduit, l'exécute (ou le compile), et restitue les résultats
ou les messages d'erreur. La figure~\ref{fig:usecase} en donne la vue « cas
d'utilisation ».

\begin{figure}[ht]
\centering
\begin{tikzpicture}[font=\small,
  uc/.style={draw,ellipse,align=center,minimum width=32mm,inner sep=2pt}]
  \node (act) at (0,0) {\begin{tikzpicture}
      \draw (0,0) circle (2.2mm);
      \draw (0,-2.2mm) -- (0,-9mm);
      \draw (-3mm,-5mm) -- (3mm,-5mm);
      \draw (0,-9mm) -- (-3mm,-14mm);
      \draw (0,-9mm) -- (3mm,-14mm);
    \end{tikzpicture}};
  \node[below=8mm of act,font=\footnotesize] {Utilisateur};
  \node[uc] (u1) at (4.5,1.6) {Exécuter un\\fichier \code{.ilang}};
  \node[uc] (u2) at (4.5,0)   {Saisir en mode\\interactif};
  \node[uc] (u3) at (4.5,-1.6){Générer\\l'assembleur};
  \draw (act) -- (u1);
  \draw (act) -- (u2);
  \draw (act) -- (u3);
  \begin{scope}[on background layer]
    \node[draw,dashed,fit=(u1)(u2)(u3),inner sep=4mm,label=above:{\footnotesize Interpréteur \ilang}] {};
  \end{scope}
\end{tikzpicture}
\caption{Diagramme des cas d'utilisation.}
\label{fig:usecase}
\end{figure}

\paragraph{Un programme \ilang{} complet.} Pour fixer les idées, voici un
programme exerçant la plupart des constructions du langage, à savoir une fonction
récursive, une conditionnelle, une boucle et un affichage :

\begin{lstlisting}[style=ilang,caption={Exemple : factorielle récursive et boucle d'affichage.}]
fonction factorielle(n) {
    si (n <= 1) {
        retourner(1)
    }
    retourner(n * factorielle(n - 1))   // appel recursif
}

pour (i = 1; i < 6; i = i + 1) {
    afficher(factorielle(i))            // affiche 1, 2, 6, 24, 120
}
\end{lstlisting}

Ce programme combine déclaration de fonction, récursion, structure
conditionnelle, boucle \code{pour} et entrée-sortie : il servira de fil
conducteur dans la suite du rapport.

\section{Conception architecturale}

\subsection{La chaîne de traitement}

L'architecture suit le modèle classique d'un \emph{pipeline} de compilation, où la
sortie de chaque étage est l'entrée du suivant. La
figure~\ref{fig:pipeline} en donne la vue d'ensemble.

\begin{figure}[ht]
\centering
\begin{tikzpicture}[
  node distance=8mm and 12mm,
  box/.style={draw,rounded corners,minimum height=9mm,minimum width=20mm,align=center,font=\small},
  data/.style={font=\small\itshape,align=center},
  arr/.style={-{Stealth}, very thick}
]
\node[data] (src) {source\\\code{.ilang}};
\node[box,right=of src] (lex) {Flex\\(lexer)};
\node[data,right=of lex] (tok) {tokens};
\node[box,right=of tok] (par) {Bison\\(parseur)};
\node[box,below=of par] (opt) {\code{optim.c}\\(AST $\to$ AST)};
\node[box,left=of opt] (eval) {\code{eval.c}\\(évaluateur)};
\node[data,left=of eval] (res) {résultat};

\draw[arr] (src) -- (lex);
\draw[arr] (lex) -- (tok);
\draw[arr] (tok) -- (par);
\draw[arr] (par) -- node[right,font=\scriptsize]{AST} (opt);
\draw[arr] (opt) -- node[above,font=\scriptsize]{AST optimisé} (eval);
\draw[arr] (eval) -- (res);

\node[box,below=14mm of eval,fill=black!5] (cg) {\code{codegen.c}\\(back-end asm)};
\draw[arr,dashed] (opt.south) |- (cg.east) node[pos=0.25,right,font=\scriptsize]{extension};
\node[data,left=of cg] (asm) {\code{.s} i386};
\draw[arr,dashed] (cg) -- (asm);
\end{tikzpicture}
\caption{Chaîne de traitement d'\ilang. Le front-end (Flex, Bison) est commun ;
seul le back-end diffère entre l'interprétation et la génération de code.}
\label{fig:pipeline}
\end{figure}

\subsection{Patrons de conception (\emph{patterns})}

Deux patrons structurent le logiciel.

\begin{itemize}
  \item \textbf{Pipeline} (chaîne de responsabilité) : chaque module transforme
        une représentation en une autre (texte $\to$ tokens $\to$ AST $\to$ AST
        optimisé $\to$ valeur). Cette séparation rend chaque étage testable et
        remplaçable indépendamment.
  \item \textbf{Parcours d'arbre type \emph{visiteur}} : l'évaluateur, l'optimiseur
        et le générateur de code sont trois « visiteurs » distincts du même AST.
        En C, sans classes, le patron se concrétise par une fonction récursive qui
        distribue (\code{switch}) sur le type du n\oe ud. C'est cette uniformité
        qui a rendu l'ajout du back-end assembleur peu coûteux : un troisième
        visiteur, sans toucher aux deux autres.
\end{itemize}

\subsection{Choix du langage d'implémentation}

L'interpréteur est écrit en \textbf{C}. Ce choix n'est pas neutre et se justifie :

\begin{enumerate}
  \item Flex et Bison génèrent du code C ; écrire le reste en C évite toute
        couche d'interfaçage entre langages.
  \item Le C impose de gérer explicitement la mémoire et les structures de
        données, ce qui est exactement l'objet du cours (pile d'appels, n\oe uds
        alloués dynamiquement).
  \item Le back-end assembleur produit du code 32~bits lié à la \emph{libc} C ;
        rester en C garantit une intégration naturelle.
\end{enumerate}

\subsection{Justification des principales décisions de conception}

Plusieurs décisions de conception méritent d'être explicitées, car à chaque fois
une autre solution était envisageable et le choix retenu résulte d'un arbitrage
que je veux assumer.

La première décision est de construire un arbre syntaxique abstrait plutôt que
d'évaluer le programme directement pendant l'analyse syntaxique. Cette seconde
option, appelée évaluation dirigée par la syntaxe, est parfaitement possible pour
une calculatrice : chaque règle de la grammaire calcule immédiatement sa valeur.
Je l'ai écartée pour une raison de fond. Évaluer pendant l'analyse couplerait
définitivement la syntaxe et la sémantique, interdirait toute passe
d'optimisation, qui suppose de disposer du programme entier avant de l'exécuter,
et rendrait impossible la réutilisation du même travail d'analyse pour produire du
code. En matérialisant le programme sous forme d'arbre, je me donne au contraire
la liberté de le parcourir autant de fois que nécessaire et de plusieurs manières.
C'est exactement cette liberté qui m'a permis, plus tard, d'ajouter un optimiseur
puis un générateur de code sans rien changer à l'analyse.

La deuxième décision concerne la table des symboles. J'aurais pu me contenter
d'une unique table globale, ce qui aurait été plus simple, mais cela aurait rendu
impossible la gestion correcte de la portée locale et de la récursion : deux
appels simultanés de la même fonction auraient partagé les mêmes variables. À
l'inverse, une structure trop élaborée aurait été disproportionnée. La pile
d'environnements, où chaque appel dispose de son propre cadre, est la structure
qui correspond exactement au besoin, et elle a en outre le mérite de reproduire
fidèlement la pile d'appels matérielle étudiée au chapitre~2. Le choix d'une table
de hachage à l'intérieur de chaque cadre, plutôt qu'une simple liste chaînée, vise
la rapidité de recherche : l'accès à une variable est l'opération la plus fréquente
de tout l'évaluateur, et il doit rester proche du temps constant.

La troisième décision, introduite avec l'extension des flottants, est de
représenter une valeur par une union étiquetée distinguant entier et flottant,
plutôt que de tout convertir en flottant. Cette dernière solution aurait été plus
courte à écrire, mais elle aurait silencieusement détruit la sémantique entière
d'origine, en particulier la division entière et le modulo dont dépend le calcul
du PGCD. L'union étiquetée coûte quelques lignes de plus à chaque opération, mais
elle préserve exactement le comportement attendu, ce qui était une contrainte non
négociable de non-régression.

La quatrième décision concerne la gestion des ruptures de flot, c'est-à-dire le
retour de fonction et les instructions \code{arrêter} et \code{continuer}. J'aurais
pu employer les sauts non locaux du C, à savoir \code{setjmp} et \code{longjmp},
ou faire remonter un code de statut par chaque appel de la fonction d'évaluation.
J'ai préféré des indicateurs globaux que les blocs et les boucles consultent. Ce
choix privilégie la lisibilité : le mécanisme tient en quelques lignes et reste
facile à suivre, ce qui compte davantage qu'une élégance théorique dans un
interpréteur mono-thread. J'assume que cette solution introduit un état global,
et j'en discute les limites au chapitre~\ref{chap:discussion}.

\section{Méthodologie et planification}

\subsection{Méthodologie retenue}

Le développement a suivi une démarche \textbf{incrémentale et pilotée par les
tests}. Plutôt que de tout écrire avant d'exécuter quoi que ce soit, le c\oe ur
(entiers seulement) a d'abord été rendu fonctionnel de bout en bout, validé sur
les programmes de démonstration, puis enrichi par étapes (optimisations, puis
extensions, puis back-end). Chaque ajout était suivi d'une vérification de
non-régression. Cette stratégie « marcher avant de courir » s'est révélée
payante : le passage délicat des entiers aux flottants, par exemple, aurait été
bien plus risqué sur une base non encore éprouvée.

\subsection{Identification des tâches et planning}

Le projet se décompose en huit grandes tâches, présentées sous forme de diagramme
de Gantt (figure~\ref{fig:gantt}). Chacune est développée dans la suite du
rapport.

\begin{figure}[ht]
\centering
\resizebox{\textwidth}{!}{%
\begin{ganttchart}[
  vgrid, hgrid,
  bar/.append style={fill=black!25},
  bar incomplete/.append style={fill=black!10},
  group/.append style={fill=black!55},
  title label font=\scriptsize,
  bar label font=\small,
  group label font=\small\bfseries,
  milestone label font=\scriptsize\itshape
]{1}{12}
  \gantttitle{Déroulement du projet (unités de temps relatives)}{12} \\
  \gantttitlelist{1,...,12}{1} \\
  \ganttgroup{Conception}{1}{2} \\
  \ganttbar{Grammaire \& tokens}{1}{2} \\
  \ganttgroup{C\oe ur de l'interpréteur}{2}{6} \\
  \ganttbar{Lexer (Flex)}{2}{3} \\
  \ganttbar{Parseur \& AST (Bison)}{3}{4} \\
  \ganttbar{Évaluateur}{4}{5} \\
  \ganttbar{Table des symboles}{5}{6} \\
  \ganttbar{Optimisations}{6}{6} \\
  \ganttgroup{Tests \& débogage}{5}{8} \\
  \ganttbar{Programmes de démo}{5}{6} \\
  \ganttbar{Gestion des erreurs}{6}{7} \\
  \ganttbar{Vérification (Valgrind)}{7}{8} \\
  \ganttgroup{Extensions}{8}{10} \\
  \ganttbar{Flottants}{8}{9} \\
  \ganttbar{arrêter / continuer}{9}{9} \\
  \ganttbar{Back-end assembleur}{9}{10} \\
  \ganttgroup{Rédaction}{10}{12} \\
  \ganttbar{Journal \& rapport}{10}{12} \\
  \ganttmilestone{Rendu}{12}
\end{ganttchart}%
}
\caption{Diagramme de Gantt du déroulement du projet.}
\label{fig:gantt}
\end{figure}

\section{Conception détaillée}

\subsection{La grammaire (BNF)}

Le c\oe ur de la conception syntaxique est la grammaire, donnée ci-dessous en
notation BNF. C'est elle que Bison transforme en analyseur LALR(1).

\begin{lstlisting}[style=gram,caption={Grammaire BNF d'\ilang.},label={lst:bnf}]
programme   ::= liste_instructions

instruction ::= affectation osemi
              | si_alors | tant_que | pour
              | declaration_fonction
              | appel_fonction osemi
              | AFFICHER '(' expression ')' osemi
              | RETOURNER '(' expression ')' osemi
              | ARRETER osemi | CONTINUER osemi

osemi       ::= /* vide */ | ';'      /* point-virgule optionnel */

affectation ::= IDENTIFIANT '=' expression

expression  ::= expression '+' expression  | expression '-' expression
              | expression '*' expression  | expression '/' expression
              | expression '%' expression
              | expression '==' expression | expression '!=' expression
              | expression '<'  expression | expression '>'  expression
              | expression '<=' expression | expression '>=' expression
              | '-' expression %prec UMINUS
              | '(' expression ')'
              | NOMBRE | FLOTTANT | IDENTIFIANT
              | appel_fonction | LIRE '(' ')'

si_alors    ::= SI '(' expression ')' bloc
              | SI '(' expression ')' bloc SINON bloc
tant_que    ::= TANTQUE '(' expression ')' bloc
pour        ::= POUR '(' affectation ';' expression ';' affectation ')' bloc
bloc        ::= '{' liste_instructions '}'

declaration_fonction ::= FONCTION IDENTIFIANT '(' liste_params ')' bloc
appel_fonction       ::= IDENTIFIANT '(' liste_args ')'
\end{lstlisting}

La priorité des opérateurs n'est pas exprimée dans la BNF elle-même (qui est
volontairement ambiguë sur ce point) mais par des directives de priorité, de la
plus faible à la plus forte : comparaisons, puis \code{+ -}, puis \code{* / \%},
puis le moins unaire. Ce mécanisme, propre à Bison, évite de devoir réécrire la
grammaire en cascade de non-terminaux (\emph{terme}, \emph{facteur}, etc.).

Un point mérite une explication, car il révèle une tension interne au cahier des
charges. La grammaire BNF de référence fait figurer un point-virgule à la fin de
l'affectation, de l'appel, de l'\code{afficher} et du \code{retourner}. Or aucun
des programmes de démonstration fournis n'utilise de point-virgule, sauf à
l'intérieur de l'en-tête de la boucle \code{pour}. Plutôt que de trancher
arbitrairement en faveur de l'une des deux sources, j'ai choisi de rendre le
point-virgule \emph{optionnel} en fin d'instruction, au moyen du non-terminal
\code{osemi} qui se réduit soit sur un point-virgule, soit sur le vide. Ainsi les
deux styles d'écriture sont acceptés : les programmes de démonstration sans
point-virgule comme les programmes rédigés à la lettre de la BNF. Cette
souplesse n'introduit aucun conflit dans la grammaire, car le point-virgule
n'apparaît comme séparateur obligatoire que dans le \code{pour}, un contexte
distinct.

\subsection{Structure de l'arbre syntaxique abstrait}

Chaque n\oe ud de l'AST est une structure C uniforme : un champ \code{type}
indiquant la nature du n\oe ud, et un ensemble de champs dont seuls certains sont
significatifs selon le type. La figure~\ref{fig:ast-uml} en donne une vue de type
UML, et la figure~\ref{fig:ast-ex} montre l'arbre obtenu pour une expression.

\begin{figure}[ht]
\centering
\begin{tikzpicture}[font=\footnotesize]
\node[draw,rectangle split,rectangle split parts=2,align=left] (n) {
  \textbf{Node}
  \nodepart{second}
  \begin{tabular}{l}
  + type : NodeType \\
  + num : Value \quad\textit{// NODE\_NUM} \\
  + name : char* \quad\textit{// VAR, ASSIGN, FUNC\_*} \\
  + op : int \quad\textit{// BINOP, UNOP} \\
  + left, right, cond : Node* \\
  + body, elsebody : Node* \\
  + args : Node** ; nargs : int \\
  + line : int \\
  \end{tabular}
};
\node[draw,rectangle split,rectangle split parts=2,align=left,right=18mm of n] (v) {
  \textbf{Value}
  \nodepart{second}
  \begin{tabular}{l}
  + type : \{VAL\_INT, VAL\_FLOAT\} \\
  + ival : long \\
  + fval : double \\
  \end{tabular}
};
\draw[-{Stealth}] (n) -- (v) node[midway,above,font=\scriptsize]{num};
\end{tikzpicture}
\caption{Structure des n\oe uds de l'AST (vue UML simplifiée).}
\label{fig:ast-uml}
\end{figure}

\begin{figure}[ht]
\centering
\begin{tikzpicture}[
  level distance=11mm, sibling distance=20mm, font=\small,
  every node/.style={draw,circle,inner sep=1.5pt,minimum size=7mm}
]
\node {\code{+}}
  child {node {\code{*}}
    child {node {2}}
    child {node {3}}
  }
  child {node {4}};
\end{tikzpicture}
\caption{AST de l'expression \code{2 * 3 + 4}. La priorité de \code{*} sur
\code{+} se traduit par la position des n\oe uds dans l'arbre.}
\label{fig:ast-ex}
\end{figure}

Le langage compte quatorze types de n\oe uds, recensés dans le
tableau~\ref{tab:nodetypes} avec les champs significatifs de chacun. Cette
uniformité (une seule structure \code{Node}, un champ \code{type}) est ce qui rend
les trois visiteurs, à savoir l'évaluateur, l'optimiseur et le générateur de
code, aussi réguliers à écrire.

\begin{table}[ht]
\centering
\caption{Types de n\oe uds de l'AST et champs significatifs.}
\label{tab:nodetypes}
\small
\begin{tabular}{@{}lll@{}}
\toprule
\textbf{Type} & \textbf{Rôle} & \textbf{Champs utilisés} \\
\midrule
\code{NODE\_NUM}       & littéral numérique      & \code{num} \\
\code{NODE\_VAR}       & référence de variable   & \code{name} \\
\code{NODE\_BINOP}     & opération binaire       & \code{op, left, right} \\
\code{NODE\_UNOP}      & moins unaire            & \code{op, left} \\
\code{NODE\_ASSIGN}    & affectation             & \code{name, left} \\
\code{NODE\_IF}        & conditionnelle          & \code{cond, body, elsebody} \\
\code{NODE\_WHILE}     & boucle tant que         & \code{cond, body} \\
\code{NODE\_FOR}       & boucle pour             & \code{left, cond, right, body} \\
\code{NODE\_FUNC\_DECL}& déclaration de fonction & \code{name, args, nargs, body} \\
\code{NODE\_FUNC\_CALL}& appel de fonction       & \code{name, args, nargs} \\
\code{NODE\_PRINT}     & \code{afficher}         & \code{left} \\
\code{NODE\_READ}      & \code{lire}             & aucun \\
\code{NODE\_RETURN}    & \code{retourner}        & \code{left} \\
\code{NODE\_BREAK} / \code{NODE\_CONTINUE} & ruptures de boucle & aucun \\
\code{NODE\_BLOCK}     & suite d'instructions    & \code{args, nargs} \\
\bottomrule
\end{tabular}
\end{table}

\subsection{La table des symboles}

La gestion de la portée repose sur une \textbf{pile d'environnements} (\emph{frames}),
chaque environnement étant lui-même une table de hachage. Un environnement est
empilé à l'entrée d'une fonction et dépilé à sa sortie, à l'image de la pile
d'appels matérialisée par le registre \code{\%ebp} en assembleur i386 (chapitre~2
du cours). La figure~\ref{fig:frames} illustre l'état de la pile pendant un appel
récursif.

\begin{figure}[ht]
\centering
\begin{tikzpicture}[font=\footnotesize, frame/.style={draw,minimum width=40mm,minimum height=8mm,align=center}]
  \node[frame] (g) {environnement \textbf{global} \\ \code{i = 3}};
  \node[frame,above=2mm of g] (f1) {frame \code{fib(3)} \\ \code{n = 3}};
  \node[frame,above=2mm of f1] (f2) {frame \code{fib(2)} \\ \code{n = 2}};
  \node[frame,above=2mm of f2,fill=black!8] (f3) {frame \code{fib(1)} \quad$\leftarrow$ sommet \\ \code{n = 1}};
  \draw[-{Stealth}] (f3.east) to[bend left=30] node[right,font=\scriptsize]{parent} (f2.east);
  \draw[-{Stealth}] (f2.east) to[bend left=30] (f1.east);
  \draw[-{Stealth}] (f1.east) to[bend left=30] (g.east);
\end{tikzpicture}
\caption{Pile d'environnements pendant l'évaluation de \code{fib(3)}. Le lien
\emph{parent} sert à restaurer le frame courant en sortie de fonction ; la
résolution des variables, elle, est lexicale (frame courant puis global).}
\label{fig:frames}
\end{figure}

\subsection{Algorithmes clés}

\paragraph{Évaluation.} L'évaluateur est une fonction récursive renvoyant une
\code{Value}. Son principe, pour les n\oe uds principaux, est résumé en
pseudo-code dans le listing~\ref{lst:eval-pseudo}.

\begin{lstlisting}[style=c,caption={Principe de l'évaluateur (pseudo-code).},label={lst:eval-pseudo}]
Value eval(Node n):
    selon n.type:
      NUM       -> renvoyer n.num
      VAR       -> chercher n.name dans l'environnement (erreur si absent)
      BINOP     -> a=eval(gauche); b=eval(droite); appliquer l'operateur
      ASSIGN    -> v=eval(expr); ranger v sous n.name; renvoyer v
      IF        -> si eval(cond) vrai: eval(alors) sinon: eval(sinon)
      WHILE/FOR -> boucler tant que eval(cond), en surveillant
                   les signaux retourner / arreter / continuer
      FUNC_CALL -> evaluer les args (frame appelant), empiler un frame,
                   lier les parametres, evaluer le corps, depiler
      RETURN    -> return_value=eval(expr); lever le drapeau de retour
      BLOCK     -> evaluer chaque instruction, en s'arretant des qu'un
                   signal de rupture de flot est leve
\end{lstlisting}

\paragraph{Optimisation.} L'optimiseur est un second parcours qui \emph{réécrit}
l'arbre. Deux transformations sont appliquées : le \emph{pliage de constantes}
(un n\oe ud binaire dont les deux fils sont des littéraux est remplacé par le
littéral résultat) et l'\emph{élimination du code mort} (une conditionnelle dont
la condition est un littéral est remplacée par la seule branche atteignable). Le
détail, y compris le garde-fou contre la division par zéro, est exposé au
chapitre~\ref{chap:implementation}.

\chapter{Implémentation}
\label{chap:implementation}

\section{Environnement technique}

L'ensemble a été développé sous GNU/Linux (Ubuntu, via WSL sur Windows). La
chaîne d'outils est entièrement libre :

\begin{table}[ht]
\centering
\caption{Outils et versions employés.}
\begin{tabular}{@{}ll@{}}
\toprule
\textbf{Outil} & \textbf{Rôle} \\
\midrule
GCC & compilateur C (et assembleur \code{-m32} pour l'extension) \\
GNU Bison 3.8 & générateur d'analyseur syntaxique LALR(1) \\
Flex 2.6 & générateur d'analyseur lexical \\
GNU Make & automatisation de la compilation \\
Valgrind & détection des fuites et erreurs mémoire \\
\TeX{} Live & rédaction du présent rapport \\
\bottomrule
\end{tabular}
\end{table}

\section{Architecture logicielle}

\subsection{Organisation du code}

Le code est découpé en modules à responsabilité unique, listés dans le
tableau~\ref{tab:fichiers}. Cette séparation reflète directement le pipeline de la
figure~\ref{fig:pipeline}.

\begin{table}[ht]
\centering
\caption{Modules du projet et leur rôle.}
\label{tab:fichiers}
\begin{tabular}{@{}ll@{}}
\toprule
\textbf{Fichier} & \textbf{Responsabilité} \\
\midrule
\code{ilang.l}            & analyse lexicale (Flex) \\
\code{ilang.y}            & grammaire et construction de l'AST (Bison) \\
\code{ast.h} / \code{ast.c}     & définition et manipulation de l'AST, type \code{Value} \\
\code{symtab.h} / \code{symtab.c} & table des symboles (pile de frames + hachage) \\
\code{eval.c}             & évaluateur récursif \\
\code{optim.c}            & optimisations sur l'AST \\
\code{codegen.c}          & back-end assembleur i386 (extension) \\
\code{main.c}             & point d'entrée \\
\code{Makefile}           & compilation automatisée \\
\bottomrule
\end{tabular}
\end{table}

\subsection{Interfaces entre modules}

Les interfaces sont minimales et explicites. Flex et Bison communiquent par le
type \code{YYSTYPE} (une union) et la table des tokens engendrée
(\code{ilang.tab.h}). Bison renseigne une variable globale \code{program\_root}
pointant la racine de l'AST. \code{main.c} appelle successivement
\code{yyparse()}, \code{optimize()} puis \code{eval\_program()} (ou
\code{generate\_asm()}). L'évaluateur et le générateur de code dépendent de
\code{ast.h} ; l'évaluateur dépend en plus de \code{symtab.h}.

\section{Le c\oe ur : du texte à l'exécution}

\subsection{Analyse lexicale (Flex)}

Le lexer reconnaît les nombres, identifiants, mots-clés, opérateurs et
délimiteurs, et ignore espaces et commentaires. Sous le capot, Flex compile les
expressions régulières en un \emph{automate fini déterministe} (chapitre~9 du
cours). La figure~\ref{fig:dfa} en donne une vue schématique pour la
reconnaissance des entiers et des identifiants : depuis l'état initial, le premier
caractère lu oriente vers l'une ou l'autre des branches, et l'automate reste dans
l'état d'acceptation tant que des caractères valides se présentent (principe du
\emph{plus long préfixe}).

\begin{figure}[ht]
\centering
\begin{tikzpicture}[font=\small,>=Stealth,node distance=22mm,
  st/.style={draw,circle,minimum size=8mm},
  acc/.style={draw,circle,double,minimum size=8mm}]
  \node[st] (i) {$s_0$};
  \node[acc,above right=6mm and 22mm of i] (num) {$N$};
  \node[acc,below right=6mm and 22mm of i] (id) {$I$};
  \draw[->] (i) -- node[above left,font=\scriptsize]{\code{[0-9]}} (num);
  \draw[->] (i) -- node[below left,font=\scriptsize]{\code{[a-zA-Z\_]}} (id);
  \draw[->] (num) edge[loop right] node[font=\scriptsize]{\code{[0-9]}} ();
  \draw[->] (id) edge[loop right] node[font=\scriptsize]{\code{[a-zA-Z0-9\_]}} ();
  \node[left=4mm of i] (start) {}; \draw[->] (start) -- (i);
\end{tikzpicture}
\caption{Automate fini (simplifié) reconnaissant les entiers ($N$) et les
identifiants ($I$). Les états doublés sont acceptants.}
\label{fig:dfa}
\end{figure}

Au-delà de ce schéma général, deux points méritent l'attention.

D'abord, le mot-clé \og tant que \fg{} comporte un espace : il est capturé par un
motif dédié placé \emph{avant} la règle des identifiants. Ensuite, les flottants
sont testés avant les entiers, et acceptent les formes \code{3.14}, \code{3.} et
\code{.5} :

\begin{lstlisting}[style=gram,caption={Extrait de \code{ilang.l} : littéraux numériques.}]
({CHIFFRE}+"."{CHIFFRE}*)|("."{CHIFFRE}+)  { yylval.fnum = strtod(yytext, NULL); return FLOTTANT; }
{CHIFFRE}+   { yylval.num = strtol(yytext, NULL, 10); return NOMBRE; }
{ID}         { yylval.str = strdup(yytext); return IDENTIFIANT; }
\end{lstlisting}

Les commentaires de bloc sont gérés par un \emph{état exclusif} \code{\%x COMMENT}.
Le mécanisme mérite d'être détaillé, car il illustre la notion d'\emph{automate à
états} vue au chapitre sur les automates finis :

\begin{lstlisting}[style=gram,caption={Gestion des commentaires de bloc par état exclusif (extrait de \code{ilang.l}).}]
"/*"                { BEGIN(COMMENT); }   /* on entre dans l'etat COMMENT */
<COMMENT>"*/"       { BEGIN(INITIAL); }   /* on en sort */
<COMMENT>\n         { /* contenu ignore, mais yylineno est incremente */ }
<COMMENT>.          { /* tout autre caractere du commentaire est ignore */ }
<COMMENT><<EOF>>    { fprintf(stderr, "Erreur lexicale ligne %d : commentaire "
                        "de bloc non termine\n", yylineno); exit(1); }
\end{lstlisting}

Une fois dans l'état \code{COMMENT}, seules les règles préfixées par
\code{<COMMENT>} sont actives : le contenu est consommé caractère par caractère
sans produire de token, jusqu'à la séquence fermante \code{*/}. La règle
\code{<COMMENT><<EOF>>} attrape le cas d'un commentaire jamais refermé.

\paragraph{Exemple de flux de tokens.} Pour bien voir ce que produit le lexer,
considérons la ligne source \code{resultat = (a + b) * 2}. L'analyseur lexical la
découpe en la suite d'unités lexicales suivante (les commentaires et espaces étant
ignorés) :

\begin{table}[ht]
\centering
\caption{Flux de tokens produit pour \code{resultat = (a + b) * 2}.}
\begin{tabular}{@{}llll@{}}
\toprule
\textbf{Texte} & \textbf{Token} & \textbf{Valeur sémantique} & \textbf{Motif Flex} \\
\midrule
\code{resultat} & \code{IDENTIFIANT} & \code{"resultat"} (strdup) & \code{[a-zA-Z\_]...} \\
\code{=}        & \code{'='}          & aucune                     & \code{"="} \\
\code{(}        & \code{'('}          & aucune                     & \code{"("} \\
\code{a}        & \code{IDENTIFIANT}  & \code{"a"}                 & \code{[a-zA-Z\_]...} \\
\code{+}        & \code{'+'}          & aucune                     & \code{"+"} \\
\code{b}        & \code{IDENTIFIANT}  & \code{"b"}                 & \code{[a-zA-Z\_]...} \\
\code{)}        & \code{')'}          & aucune                     & \code{")"} \\
\code{*}        & \code{'*'}          & aucune                     & \code{"*"} \\
\code{2}        & \code{NOMBRE}       & \code{2} (strtol)          & \code{[0-9]+} \\
\bottomrule
\end{tabular}
\end{table}

C'est cette suite de tokens, et non le texte brut, que reçoit l'analyseur
syntaxique.

\subsection{Analyse syntaxique et construction de l'AST (Bison)}

La grammaire (listing~\ref{lst:bnf}) est traduite presque mot pour mot en règles
Bison, chaque règle portant une \emph{action sémantique} qui construit le n\oe ud
d'AST correspondant. Par exemple, la règle d'addition s'écrit :

\begin{lstlisting}[style=gram,caption={Action sémantique d'une règle (extrait de \code{ilang.y}).}]
expression : expression '+' expression { $$ = new_binop(OP_ADD, $1, $3); }
\end{lstlisting}

Le type des valeurs sémantiques portées par les symboles est une union, déclarée
ainsi :

\begin{lstlisting}[style=gram,caption={Union sémantique et déclaration des tokens (extrait de \code{ilang.y}).}]
%union {
    long   num;     /* valeur d'un NOMBRE  (entier 64 bits) */
    double fnum;    /* valeur d'un FLOTTANT */
    char  *str;     /* nom d'un IDENTIFIANT */
    Node  *node;    /* sous-arbre construit */
    NodeList *list; /* liste de parametres/arguments en construction */
}
%token <num>  NOMBRE
%token <fnum> FLOTTANT
%token <str>  IDENTIFIANT
\end{lstlisting}

La priorité et l'associativité des opérateurs sont déclarées en tête de fichier,
de la plus faible à la plus forte. C'est ce bloc, et non la forme de la grammaire,
qui lève l'ambiguïté de la règle « plate » des expressions :

\begin{lstlisting}[style=gram,caption={Déclarations de priorité (extrait de \code{ilang.y}).}]
%left EQ NE LT GT LE GE   /* comparaisons : priorite la plus faible */
%left '+' '-'             /* addition / soustraction */
%left '*' '/' '%'         /* multiplication / division / modulo */
%right UMINUS             /* moins unaire : priorite la plus forte */
\end{lstlisting}

Concrètement, \code{\%left '*'} placé \emph{après} \code{\%left '+'} signifie que
\code{*} lie plus fort que \code{+} : l'expression \code{2 + 3 * 4} est donc
analysée comme \code{2 + (3 * 4)}. L'associativité à gauche (\code{\%left})
impose que \code{8 - 3 - 2} se lise \code{(8 - 3) - 2}.

\paragraph{Exemple de construction d'AST.} Suivons la construction de l'arbre pour
l'affectation \code{x = 2 * 3 + 4}. À mesure que Bison réduit les règles, les
actions sémantiques appellent les constructeurs et font « remonter » des
sous-arbres via la pile sémantique :

\begin{enumerate}
  \item \code{NOMBRE(2)} se réduit en \code{expression} $\to$ \code{new\_num(2)} ;
        de même pour 3 et 4.
  \item La sous-expression \code{2 * 3} est réduite la première (priorité de
        \code{*}) $\to$ \code{new\_binop(OP\_MUL, num(2), num(3))}.
  \item Puis l'addition $\to$ \code{new\_binop(OP\_ADD, mul, num(4))}.
  \item Enfin l'affectation $\to$ \code{new\_assign("x", add)}.
\end{enumerate}

L'arbre obtenu est celui de la figure~\ref{fig:ast-assign}. On remarque que la
priorité, déclarée en quelques lignes, se matérialise dans la \emph{forme} de
l'arbre : la multiplication est plus « profonde » que l'addition, donc évaluée
avant.

\begin{figure}[ht]
\centering
\begin{tikzpicture}[level distance=11mm, sibling distance=18mm, font=\small,
  every node/.style={draw,rounded corners,inner sep=2pt,minimum size=7mm}]
\node {\code{ASSIGN x}}
  child {node {\code{+}}
    child {node {\code{*}}
      child {node {2}}
      child {node {3}}
    }
    child {node {4}}
  };
\end{tikzpicture}
\caption{AST de l'instruction \code{x = 2 * 3 + 4}.}
\label{fig:ast-assign}
\end{figure}

Les listes de paramètres et d'arguments sont accumulées dans une petite structure
jetable \code{NodeList}, dont le tableau interne est ensuite transféré dans le
n\oe ud final (\code{NODE\_FUNC\_DECL} ou \code{NODE\_FUNC\_CALL}), puis seule la
coquille est libérée, une subtilité mémoire que Valgrind a confirmée sans fuite.

\subsection{L'évaluateur}

L'évaluateur réalise le parcours décrit au listing~\ref{lst:eval-pseudo}. Les
ruptures de flot non locales (retour de fonction, \code{arrêter}, \code{continuer})
sont gérées par des drapeaux globaux que les blocs et les boucles surveillent. Le
traitement d'un appel de fonction illustre la gestion de la pile d'environnements :

\begin{lstlisting}[style=c,caption={Appel de fonction : ordre d'évaluation et pile (extrait de \code{eval.c}).}]
/* 1) Evaluer les arguments dans l'environnement de l'APPELANT. */
for (i = 0; i < node->nargs; i++) argvals[i] = eval(node->args[i]);
/* 2) Empiler un nouveau frame et y lier les parametres. */
env_push();
for (i = 0; i < decl->nargs; i++) env_set(decl->args[i]->name, argvals[i]);
/* 3) Evaluer le corps, recuperer la valeur de retour. */
return_flag = 0; eval(decl->body);
result = return_flag ? return_value : make_int(0);
/* 4) Depiler le frame. */
env_pop();
\end{lstlisting}

L'ordre est crucial : les arguments sont évalués \emph{avant} l'empilement du
nouveau frame, faute de quoi une expression comme \code{fib(n-1)} chercherait
\code{n} dans le frame vide de la fonction appelée.

Le c\oe ur du calcul arithmétique se trouve dans \code{eval\_binop}, qui applique
la règle de cohabitation entiers/flottants. En voici un extrait commenté :

\begin{lstlisting}[style=c,caption={Cohabitation entiers/flottants et garde-fou (extrait de \code{eval.c}).}]
Value a = eval(node->left);
Value b = eval(node->right);
int both_int = (a.type == VAL_INT && b.type == VAL_INT);
switch (node->op) {
case OP_ADD:                       /* si deux entiers -> entier */
    return both_int ? make_int(a.ival + b.ival)
                    : make_float(as_double(a) + as_double(b)); /* sinon flottant */
case OP_DIV:
    if (both_int) {
        if (b.ival == 0) {         /* division entiere par zero : */
            fprintf(stderr, "Erreur ligne %d : division par zero\n", node->line);
            return make_int(0);    /* message emis, le programme continue */
        }
        return make_int(a.ival / b.ival);  /* division ENTIERE preservee */
    }
    /* ... cas flottant ... */
}
\end{lstlisting}

\paragraph{Exemple : trace d'exécution de \code{fib(3)}.} Pour rendre concret le
mécanisme d'appel récursif et de pile d'environnements, déroulons l'évaluation de
\code{fib(3)} avec la définition usuelle. Le tableau~\ref{tab:trace-fib} montre
l'enchaînement des appels, l'état du sommet de pile et la valeur retournée.

\begin{table}[ht]
\centering
\caption{Trace simplifiée de l'évaluation de \code{fib(3)}.}
\label{tab:trace-fib}
\small
\begin{tabular}{@{}llll@{}}
\toprule
\textbf{Appel} & \textbf{\code{n}} & \textbf{Action} & \textbf{Retour} \\
\midrule
\code{fib(3)} & 3 & $n \geq 2$ : calcule \code{fib(2)+fib(1)} & 2 \\
\quad\code{fib(2)} & 2 & $n \geq 2$ : calcule \code{fib(1)+fib(0)} & 1 \\
\quad\quad\code{fib(1)} & 1 & $n < 2$ : \code{retourner(1)} & 1 \\
\quad\quad\code{fib(0)} & 0 & $n < 2$ : \code{retourner(0)} & 0 \\
\quad\code{fib(1)} & 1 & $n < 2$ : \code{retourner(1)} & 1 \\
\bottomrule
\end{tabular}
\end{table}

Chaque ligne indentée correspond à un \code{env\_push} (nouveau frame avec son
propre \code{n}) suivi, au retour, d'un \code{env\_pop}. Le drapeau de retour
(\code{return\_flag}) est levé par \code{retourner} et « consommé » par la boucle
d'appel, qui récupère alors \code{return\_value}. Le résultat final, \code{fib(3) =
2}, est conforme à l'exécution réelle (cf. chapitre~\ref{chap:validation}).

\section{Aspects techniques notables}

\subsection{Le type \code{Value} : faire cohabiter entiers et flottants}

L'extension des flottants a transformé le type des valeurs manipulées. Plutôt que
de tout passer en \code{double} (ce qui aurait cassé la division entière et le
modulo), une \emph{union étiquetée} distingue les deux natures :

\begin{lstlisting}[style=c,caption={Le type \code{Value} (extrait de \code{ast.h}).}]
typedef enum { VAL_INT, VAL_FLOAT } ValType;
typedef struct Value { ValType type; long ival; double fval; } Value;
\end{lstlisting}

La règle sémantique est simple : une opération entre deux entiers reste entière ;
dès qu'un flottant intervient, le calcul bascule en \code{double}. Ainsi
\code{10 / 4} vaut \code{2} (sémantique d'origine préservée) tandis que
\code{10.0 / 4} vaut \code{2.5}.

\subsection{La portée lexicale}

Chaque environnement est une table de hachage de 211 alvéoles (un nombre premier,
pour mieux répartir les collisions), indexée par la fonction \code{djb2} de Dan
Bernstein :

\begin{lstlisting}[style=c,caption={Fonction de hachage djb2 (extrait de \code{symtab.c}).}]
static unsigned int hash(const char *s) {
    unsigned int h = 5381;
    int c;
    while ((c = (unsigned char)*s++) != 0)
        h = ((h << 5) + h) + c;   /* h * 33 + c */
    return h % HASH_SIZE;
}
\end{lstlisting}

La résolution d'une variable se fait dans le frame courant puis, en repli, dans
l'environnement global, \emph{sans} traverser les frames intermédiaires de la pile
d'appel. C'est ce qui garantit une portée \textbf{lexicale} : une fonction ne voit
jamais les variables locales de la fonction qui l'a appelée. Ce point a fait
l'objet d'une correction après une première version (involontairement) à portée
dynamique, détaillée au chapitre~\ref{chap:discussion}.

\paragraph{Exemple : la table des symboles pendant \code{pgcd(48, 18)}.} Lors de
l'appel à la fonction \code{pgcd}, un frame est empilé contenant les paramètres
\code{a} et \code{b}, auxquels s'ajoute la variable locale \code{temp} créée par
la première affectation. Le tableau~\ref{tab:pgcd-frame} montre l'évolution du
contenu de ce frame au fil des itérations de la boucle \code{tant que}.

\begin{table}[ht]
\centering
\caption{Évolution du frame de \code{pgcd} (algorithme d'Euclide sur 48 et 18).}
\label{tab:pgcd-frame}
\begin{tabular}{@{}cccc@{}}
\toprule
\textbf{Itération} & \code{a} & \code{b} & \code{temp} \\
\midrule
entrée   & 48 & 18 & (non créé) \\
1        & 18 & 12 & 18 \\
2        & 12 & 6  & 12 \\
3        & 6  & 0  & 6 \\
\midrule
\multicolumn{4}{c}{$b = 0$ : sortie de boucle, \code{retourner(a)} $\Rightarrow$ \textbf{6}} \\
\bottomrule
\end{tabular}
\end{table}

\subsection{Localisation des erreurs à la ligne}

Tous les messages d'erreur indiquent la ligne fautive, au format
\code{Erreur ligne N : ...}. Le choix de conception décisif est de \emph{figer le
numéro de ligne dans chaque n\oe ud au moment du parsing}
(\code{node->line = yylineno} dans le constructeur), et non à l'évaluation. En
effet, au moment où l'évaluation commence, l'analyse lexicale est terminée :
\code{yylineno} vaut alors la \emph{dernière} ligne du fichier. L'utiliser dans
l'évaluateur ferait pointer toutes les erreurs sur la fin du fichier. En gravant
la ligne dans le n\oe ud à sa création, chaque diagnostic d'exécution désigne le
bon emplacement, comme le confirme le test du « mauvais nombre d'arguments », qui
signale correctement la ligne~5, celle de l'appel fautif, et non la fin du
fichier.

\subsection{Les optimisations et le garde-fou de la division par zéro}

Le pliage de constantes calcule à la compilation les sous-expressions
entièrement constantes. Un piège guette : \code{10 / 0} est, lui aussi, composé de
deux littéraux. Le plier provoquerait soit un \code{SIGFPE} pendant la
compilation, soit un court-circuitage du message d'erreur attendu à l'exécution.
Une fonction \code{can\_fold} interdit donc de plier toute division ou modulo qui
échouerait ; ces cas restent dans l'arbre et c'est l'évaluateur qui émet le
diagnostic.

\begin{lstlisting}[style=c,caption={Garde-fou du pliage (extrait de \code{optim.c}).}]
static int can_fold(int op, Value a, Value b) {
    int both_int = (a.type == VAL_INT && b.type == VAL_INT);
    if (op == OP_DIV) return both_int ? (b.ival != 0) : (as_double(b) != 0.0);
    if (op == OP_MOD) return both_int && (b.ival != 0); /* modulo : entiers, != 0 */
    return 1;   /* toutes les autres operations sont pliables */
}
\end{lstlisting}

\paragraph{Exemple : avant/après pliage.} Sur l'expression \code{2 + 3 * 4},
l'optimiseur reconnaît que \code{3 * 4} a deux fils littéraux et le remplace par
\code{12}, puis que \code{2 + 12} est à son tour pliable, et le remplace par
\code{14}. L'arbre de gauche (figure~\ref{fig:fold}) est ainsi réduit à une
unique feuille. À l'exécution, \code{afficher(2 + 3 * 4)} n'effectue donc
\emph{aucun} calcul : la valeur 14 est déjà dans l'arbre.

\begin{figure}[ht]
\centering
\begin{tikzpicture}[level distance=10mm, sibling distance=15mm, font=\small,
  every node/.style={draw,rounded corners,inner sep=2pt,minimum size=6mm}]
\node {\code{+}}
  child {node {2}}
  child {node {\code{*}}
    child {node {3}}
    child {node {4}}
  };
\end{tikzpicture}
\hspace{1.5cm}
\raisebox{1.2cm}{$\xrightarrow{\text{ pliage }}$}
\hspace{1.5cm}
\raisebox{1.0cm}{%
\begin{tikzpicture}[font=\small, every node/.style={draw,rounded corners,inner sep=2pt,minimum size=6mm}]
\node {14};
\end{tikzpicture}}
\caption{Pliage de constantes : l'arbre de \code{2 + 3 * 4} est réduit à la
feuille \code{14}.}
\label{fig:fold}
\end{figure}

De même, l'élimination du code mort transforme \code{si (0) \{ A \} sinon \{ B \}}
en le seul bloc \code{B}, et \code{si (1) \{ A \}} en le seul bloc \code{A}, la
condition constante étant évaluée une fois pour toutes à la compilation.

\subsection{Gestion de la mémoire}

L'AST est entièrement alloué dynamiquement, puis libéré récursivement par
\code{free\_node}. Le transfert du tableau d'arguments depuis la \code{NodeList}
vers le n\oe ud (sans le dupliquer ni le libérer deux fois) a demandé une
attention particulière. La vérification Valgrind confirme l'absence de fuite (voir
chapitre~\ref{chap:validation}), après ajout d'un appel à \code{yylex\_destroy()}
pour libérer le tampon interne de Flex.

\subsection{Sûreté et robustesse}

Bien qu'\ilang{} ne soit pas exposé à un contexte adverse, plusieurs propriétés de
\emph{sûreté} ont été soignées. La mémoire est intégralement maîtrisée :
allocations vérifiées (tout \code{malloc}/\code{realloc} dont l'échec est fatal
est testé), libération récursive de l'AST, et absence de fuite confirmée par
Valgrind. Aucune entrée utilisateur n'est copiée dans un tampon de taille fixe :
les identifiants sont dupliqués par \code{strdup} (allocation dimensionnée), ce
qui exclut tout débordement de tampon à l'analyse. Les conditions d'erreur
(variable ou fonction indéfinie, mauvais nombre d'arguments, division par zéro,
rupture de flot hors boucle, caractère ou syntaxe invalides) sont toutes
interceptées et signalées plutôt que de produire un comportement indéfini
silencieux. La seule limite de sûreté identifiée est le débordement de la pile de
l'hôte en cas de récursion très profonde (cf. chapitre~\ref{chap:validation}),
inhérent à l'évaluation récursive.

\subsection{Le back-end assembleur (extension hors-périmètre)}

Le module \code{codegen.c} parcourt le même AST mais émet du code assembleur i386
(syntaxe AT\&T) au lieu d'évaluer. Il applique la convention d'appel cdecl :
arguments empilés de droite à gauche, résultat dans \code{\%eax}, cadre de pile
géré par \code{\%ebp}/\code{\%esp}. Les variables du programme principal deviennent
des globales en \code{.bss}, les paramètres et locales des fonctions sont rangés
en déplacement par rapport à \code{\%ebp}. Les expressions sont compilées en
machine à pile. Ce module, entièrement optionnel, n'altère en rien
l'interpréteur ; il est décrit en détail au chapitre~\ref{chap:discussion} et son
code figure en annexe~\ref{annexe:code}.

La traduction d'une opération binaire illustre la stratégie de machine à pile :

\begin{lstlisting}[style=c,caption={Génération d'un binop en machine à pile (extrait de \code{codegen.c}).}]
gen_expr(out, n->left);                 /* operande gauche -> %eax */
fprintf(out, "\tpushl %%eax\n");        /* sauvegarde sur la pile */
gen_expr(out, n->right);                /* operande droit  -> %eax */
fprintf(out, "\tmovl %%eax, %%ecx\n");  /* droit  -> %ecx */
fprintf(out, "\tpopl %%eax\n");         /* gauche -> %eax (restaure) */
switch (n->op) {
case OP_ADD: fprintf(out, "\taddl %%ecx, %%eax\n"); break;  /* eax = eax + ecx */
case OP_MUL: fprintf(out, "\timull %%ecx, %%eax\n"); break;
/* division : cltd etend le signe, idivl place le quotient dans %eax */
case OP_DIV: fprintf(out, "\tcltd\n\tidivl %%ecx\n"); break;
}
\end{lstlisting}

\paragraph{Exemple : assembleur produit pour une fonction.} Pour la fonction
\code{fonction carre(x) \{ retourner(x * x) \}}, le générateur produit le code du
listing~\ref{lst:asm-carre}, ici annoté. On y reconnaît le prologue cdecl, l'accès
au paramètre par déplacement positif depuis \code{\%ebp}, la machine à pile pour
le produit, et l'épilogue.

\begin{lstlisting}[style=asm,caption={Assembleur généré pour \code{carre} (annoté).},label={lst:asm-carre}]
il_carre:
    pushl %ebp            # prologue : sauvegarde l'ancien base pointer
    movl %esp, %ebp       # etablit le nouveau cadre de pile
    movl 8(%ebp), %eax    # charge le parametre x (1er arg en +8)
    pushl %eax            # empile x (operande gauche)
    movl 8(%ebp), %eax    # recharge x (operande droit)
    movl %eax, %ecx       # x -> %ecx
    popl %eax             # x -> %eax
    imull %ecx, %eax      # %eax = x * x
    leave                 # epilogue : restaure %esp et %ebp (retourner)
    ret                   # retour, resultat dans %eax
\end{lstlisting}

Côté appelant, \code{afficher(carre(7))} empile la constante 7, appelle
\code{il\_carre}, nettoie la pile (\code{addl \$4, \%esp}), puis transmet le
résultat (\code{\%eax}) à \code{printf} via le format \code{"\%d\textbackslash n"}.
Cette correspondance directe entre les concepts de l'évaluateur (pile d'appels,
liaison des paramètres) et leur réalisation matérielle (cadre \code{\%ebp},
déplacements) est précisément ce que le chapitre~2 du cours cherche à faire
comprendre.

\paragraph{Traduction des boucles.} Les structures de contrôle deviennent des
\emph{labels} et des \emph{sauts}. Une boucle \code{tant que (c) \{ corps \}} est
encadrée par deux labels : un label de tête (avant le test) et un label de fin.
La condition échouée saute vers la fin ; le corps se termine par un saut
inconditionnel vers la tête. Les instructions \code{arrêter} et \code{continuer}
se traduisent alors trivialement par un \code{jmp} vers, respectivement, le label
de fin et le label de tête (ou d'incrément pour une boucle \code{pour}). Le
générateur maintient à cet effet une pile de labels de boucles, exactement comme
l'évaluateur maintient un signal de rupture. Schématiquement :

\begin{lstlisting}[style=asm,caption={Schéma de traduction d'une boucle \code{tant que}.}]
.Ltete:                  # etiquette de tete
    < code de la condition c, resultat dans %eax >
    cmpl $0, %eax
    jz .Lfin             # si c est faux -> sortir
    < code du corps >    # 'arreter' = jmp .Lfin ; 'continuer' = jmp .Ltete
    jmp .Ltete           # boucler
.Lfin:                   # etiquette de fin
\end{lstlisting}

\subsection{Compilation automatisée}

Le \code{Makefile} sépare deux jeux d'options : le code du projet est compilé avec
\code{-Wall -Wextra} (zéro avertissement exigé), tandis que le code engendré par
Flex et Bison est compilé sans \code{-Wextra}, ces fichiers machine-générés
pouvant contenir des constructions qui déclencheraient des avertissements non
imputables au projet. Les dépendances entre fichiers générés (\code{lex.yy.c}
dépend de \code{ilang.tab.h}) sont explicitées pour que \code{make} ordonne
correctement la compilation.

\begin{lstlisting}[style=gram,caption={Séparation des jeux d'options (extrait du \code{Makefile}).}]
CFLAGS   = -Wall -Wextra -std=gnu11 -g    # notre code : tous les avertissements
GENFLAGS = -std=gnu11 -g                  # code genere : sans -Wextra

# Bison : grammaire -> parseur + en-tete des tokens
ilang.tab.c ilang.tab.h: ilang.y
	$(YACC) -d -o ilang.tab.c ilang.y
# Flex : depend de l'en-tete des tokens engendre par Bison
lex.yy.c: ilang.l ilang.tab.h
	$(LEX) -o lex.yy.c ilang.l
\end{lstlisting}

Le point d'entrée \code{main.c} se charge enfin d'analyser les arguments (chemin
du fichier, options \code{-s} et \code{--no-opt}), d'enchaîner
\code{yyparse()} $\to$ \code{optimize()} $\to$ \code{eval\_program()} (ou
\code{generate\_asm()}), puis de libérer l'AST et la table des symboles avant de
rendre la main.

\chapter{Validation et évaluation}
\label{chap:validation}

\section{Stratégie de test}

La validation combine quatre niveaux complémentaires :

\begin{description}[style=nextline,leftmargin=1em]
  \item[Tests fonctionnels.] Exécution des programmes de démonstration et
        comparaison des sorties avec les résultats attendus.
  \item[Tests de gestion d'erreurs.] Soumission de programmes fautifs pour
        vérifier que chaque catégorie d'erreur est détectée, localisée et
        signalée.
  \item[Tests de non-régression.] Après chaque évolution (optimisations,
        extensions, corrections), re-vérification que les sorties de référence
        restent identiques.
  \item[Tests de robustesse mémoire et cas limites.] Analyse Valgrind et
        campagne de cas extrêmes (fichier vide, EOF, grands nombres, récursion
        profonde, ruptures de flot mal placées).
\end{description}

\section{Protocole expérimental}

\subsection{Conditions et métriques}

Les mesures de performance ont été réalisées sur une machine équipée d'un
processeur AMD Ryzen~5 5500 (12 fils d'exécution), sous Ubuntu (WSL). Chaque temps
rapporté est le \emph{meilleur} de trois exécutions (pour limiter le bruit lié à
l'ordonnancement), mesuré par \code{/usr/bin/time}. Les métriques retenues sont :
le temps d'exécution (secondes), le résultat calculé (pour la correction), et le
verdict Valgrind (octets perdus, erreurs).

\subsection{Jeux d'essai}

Les programmes de démonstration servent de jeu d'essai fonctionnel : Fibonacci
récursif, PGCD d'Euclide, programme interactif, plus les deux programmes
illustrant les extensions (flottants, boucles avec \code{arrêter}/\code{continuer}).
Les benchmarks utilisent des programmes paramétrés (Fibonacci à différentes
tailles, boucle de 30~millions d'itérations, récursion profonde).

\subsection{Reproductibilité}

Les tests sont automatisés et rejouables. La cible \code{make test} compile puis
exécute l'ensemble des programmes de démonstration ; un script dédié soumet les
programmes fautifs et affiche, pour chacun, le message d'erreur et le code de
sortie. Ce script tient en quelques lignes et illustre bien la simplicité du
dispositif :

\begin{lstlisting}[style=gram,caption={C\oe ur du script de test des erreurs.}]
for f in variable divzero fonction arguments lexical syntaxe; do
    echo "===== $f ====="
    ./ilang "tests/err_$f.ilang"   # le programme fautif
    echo "(code de sortie : $?)"   # 1 si erreur fatale, 0 si poursuite
done
\end{lstlisting}

Cette automatisation a servi de filet de sécurité tout au long du développement.
Après chaque modification, une simple commande permettait de vérifier qu'aucune
sortie de référence n'avait changé. C'est précisément ce dispositif qui m'a permis
de garantir la non-régression lors de l'ajout des flottants, puis des deux autres
extensions, puis du back-end assembleur, sans craindre de casser silencieusement
ce qui fonctionnait déjà.

\section{Résultats fonctionnels}

\subsection{Programmes de démonstration}

Les trois programmes du cahier des charges produisent les sorties attendues.

\begin{lstlisting}[style=ilang,caption={Sortie de \code{fibonacci.ilang} (15 premiers termes).}]
0 1 1 2 3 5 8 13 21 34 55 89 144 233 377
\end{lstlisting}

\begin{lstlisting}[style=ilang,caption={Sortie de \code{pgcd.ilang} : pgcd(48, 18).}]
6
\end{lstlisting}

\begin{lstlisting}[style=ilang,caption={Sortie de \code{interactif.ilang} pour les entrées 7 puis $-5$.}]
999      999
49       5
\end{lstlisting}

Les programmes d'extension confirment le bon fonctionnement des flottants et des
ruptures de boucle :

\begin{lstlisting}[style=ilang,caption={Sortie de \code{flottants.ilang}.}]
6.28      (3.14 * 2)
3.5       (2.5 + 1)
2         (10 / 4   : division entiere preservee)
2.5       (10.0 / 4 : division reelle)
19.6349   (aire d'un disque de rayon 2.5)
1         (3.5 > 3)
\end{lstlisting}

\begin{lstlisting}[style=ilang,caption={Sortie de \code{boucles.ilang} (arrêter puis continuer puis arrêter dans un tant que).}]
0 1 2 3 4        (arret a 5)
1 3 5 7 9        (saut des pairs)
10 9 8           (arret a 7)
\end{lstlisting}

\subsection{Gestion des erreurs}

Chaque catégorie d'erreur du cahier des charges est traitée et localisée à la
ligne (tableau~\ref{tab:erreurs}). Cas notable, la division par zéro émet un
message \emph{et laisse le programme continuer}, conformément à la spécification.

\begin{table}[ht]
\centering
\caption{Gestion des erreurs : message émis et code de sortie.}
\label{tab:erreurs}
\small
\begin{tabular}{@{}lll@{}}
\toprule
\textbf{Erreur} & \textbf{Message} & \textbf{Sortie} \\
\midrule
Variable non définie & \code{Erreur ligne 1 : variable 'z' non definie} & 1 \\
Division par zéro & \code{Erreur ligne 1 : division par zero} & 0 (continue) \\
Fonction non définie & \code{Erreur ligne 1 : fonction 'inconnue' non definie} & 1 \\
Mauvais nb d'arguments & \code{... attend 2 argument(s), 1 fourni(s)} & 1 \\
Caractère invalide & \code{Erreur lexicale ligne 1 : caractere invalide '@'} & 1 \\
Erreur de syntaxe & \code{Erreur de syntaxe ligne 1 : syntax error...} & 1 \\
Rupture hors boucle & \code{Erreur ligne 2 : 'arreter' utilise hors d'une boucle} & 1 \\
\bottomrule
\end{tabular}
\end{table}

Pour illustrer concrètement le traitement des erreurs, voici les sorties réelles
produites par l'interpréteur sur une série de programmes fautifs. On notera la
localisation systématique à la ligne et, pour la division par zéro, la poursuite
de l'exécution.

\begin{lstlisting}[style=gram,caption={Sorties réelles sur des programmes fautifs.}]
$ ./ilang tests/err_variable.ilang
Erreur ligne 1 : variable 'z' non definie

$ ./ilang tests/err_divzero.ilang
Erreur ligne 1 : division par zero
0
123                  <- le programme continue apres l'erreur

$ ./ilang tests/err_fonction.ilang
Erreur ligne 1 : fonction 'inconnue' non definie

$ ./ilang tests/err_arguments.ilang
Erreur ligne 5 : fonction 'pgcd' attend 2 argument(s), 1 fourni(s)

$ ./ilang tests/err_lexical.ilang
Erreur lexicale ligne 1 : caractere invalide '@'

$ ./ilang tests/err_syntaxe.ilang
Erreur de syntaxe ligne 1 : syntax error, unexpected IDENTIFIANT, expecting '('
\end{lstlisting}

\section{Résultats de performance}

\subsection{Interprétation contre code natif}

Le tableau~\ref{tab:bench-fib} et la figure~\ref{fig:bench-fib} comparent le temps
de calcul de Fibonacci par l'interpréteur et par le binaire issu du back-end
assembleur. L'écart est spectaculaire : sur \code{fib(32)}, le code natif est de
l'ordre de \emph{plusieurs centaines de fois} plus rapide. Cela illustre
concrètement le coût du parcours d'arbre (distribution sur le type de n\oe ud,
appels de fonction C en cascade) face à du code machine compilé.

\begin{table}[ht]
\centering
\caption{Temps d'exécution de Fibonacci : interpréteur vs assembleur i386.}
\label{tab:bench-fib}
\begin{tabular}{@{}rrrr@{}}
\toprule
$n$ & résultat & interpréteur (s) & asm i386 (s) \\
\midrule
28 & 317\,811   & 0,43 & $\approx$ 0,00 \\
30 & 832\,040   & 1,13 & $\approx$ 0,00 \\
32 & 2\,178\,309 & 2,97 & 0,01 \\
\bottomrule
\end{tabular}
\end{table}

\begin{figure}[ht]
\centering
\begin{tikzpicture}
\begin{axis}[
  width=0.8\textwidth, height=6cm,
  xlabel={$n$}, ylabel={temps (s)},
  legend pos=north west,
  xtick={28,30,32},
  ymin=0, grid=both, grid style={black!10},
]
\addplot[mark=*,thick] coordinates {(28,0.43) (30,1.13) (32,2.97)};
\addplot[mark=square*,thick,densely dashed] coordinates {(28,0.002) (30,0.004) (32,0.01)};
\legend{interpréteur, asm i386}
\end{axis}
\end{tikzpicture}
\caption{Croissance du temps d'exécution de Fibonacci. La courbe de
l'interpréteur révèle bien la complexité exponentielle de l'algorithme ; le code
natif, sur la même échelle, reste collé à l'axe.}
\label{fig:bench-fib}
\end{figure}

\subsection{Effet des optimisations}

Pour isoler l'effet du pliage de constantes, un programme effectue 30~millions
d'itérations d'un corps de boucle contenant une expression constante
(\code{s = s + (2*3 + 4*5 - 6)}, qui se replie en \code{s + 20}). Le drapeau
\code{--no-opt} permet de désactiver l'optimiseur. Le pliage fait gagner environ
un tiers du temps (figure~\ref{fig:bench-opt}).

\begin{figure}[ht]
\centering
\begin{tikzpicture}
\begin{axis}[
  width=0.6\textwidth, height=5cm,
  ybar, bar width=28pt,
  symbolic x coords={avec optim., sans optim.},
  xtick=data, ymin=0,
  ylabel={temps (s)}, nodes near coords,
  grid=both, grid style={black!10},
]
\addplot coordinates {(avec optim.,4.05) (sans optim.,6.02)};
\end{axis}
\end{tikzpicture}
\caption{Temps de la boucle de 30~M d'itérations, avec et sans pliage de
constantes. Le gain ($\approx$ 33~\%) provient de la suppression du calcul
arithmétique répété à chaque tour.}
\label{fig:bench-opt}
\end{figure}

\subsection{Limite de profondeur de récursion}

Comme l'évaluateur récurse en C (un appel \ilang{} se traduit par une cascade
d'appels de \code{eval}), la profondeur de récursion est bornée par la pile de
l'hôte. La sonde du tableau~\ref{tab:rec} montre que \code{somme(5000)} aboutit,
mais que \code{somme(10000)} provoque un débordement de pile (signal SIGSEGV,
code de sortie 139). C'est une limite de fond de l'approche par parcours d'arbre,
discutée au chapitre~\ref{chap:discussion}.

\begin{table}[ht]
\centering
\caption{Profondeur de récursion : \code{somme(n)} récursive.}
\label{tab:rec}
\begin{tabular}{@{}rl@{}}
\toprule
$n$ & résultat \\
\midrule
1\,000  & 500\,500 (ok) \\
5\,000  & 12\,502\,500 (ok) \\
10\,000 & débordement de pile (SIGSEGV) \\
\bottomrule
\end{tabular}
\end{table}

\subsection{Vérification mémoire}

L'exécution sous Valgrind, après ajout de \code{yylex\_destroy()}, ne signale
\textbf{aucune fuite ni erreur} sur l'ensemble des programmes :

\begin{lstlisting}[style=gram,caption={Verdict Valgrind (programme \code{pgcd.ilang}).}]
in use at exit: 0 bytes in 0 blocks
ERROR SUMMARY: 0 errors from 0 contexts (suppressed: 0 from 0)
\end{lstlisting}

\subsection{Équivalence interpréteur / code généré}

Enfin, pour chaque programme entier, la sortie du binaire assemblé a été comparée
à celle de l'interpréteur : elles sont \textbf{identiques} (Fibonacci, PGCD,
boucles avec \code{arrêter}/\code{continuer}, et même le programme interactif qui
passe par \code{scanf}/\code{printf}). Le back-end refuse proprement les
programmes utilisant des flottants. Cette équivalence constitue une validation
croisée précieuse : deux back-ends indépendants, partant du même AST, calculent
les mêmes résultats.

\section{Étude de cas de bout en bout}

Pour conclure ce chapitre, suivons le programme \code{factorielle} (présenté au
chapitre~\ref{chap:conception}) tout au long de la chaîne, afin de relier les
résultats aux étapes décrites.

\begin{enumerate}
  \item \textbf{Analyse lexicale.} Le source est découpé en tokens :
        \code{FONCTION}, \code{IDENTIFIANT(factorielle)}, \code{'('},
        \code{IDENTIFIANT(n)}, \code{')'}, \code{'\{'}, \code{SI}, \dots{} Les
        commentaires et espaces sont éliminés.
  \item \textbf{Analyse syntaxique.} Bison réduit ces tokens selon la grammaire et
        bâtit l'AST : une \code{NODE\_FUNC\_DECL} (paramètre \code{n}, corps
        contenant un \code{NODE\_IF} et un \code{NODE\_RETURN} avec un
        \code{NODE\_BINOP} \code{*} dont le fils droit est un
        \code{NODE\_FUNC\_CALL} récursif).
  \item \textbf{Optimisation.} Le sous-arbre \code{n - 1} n'est pas constant (il
        dépend de \code{n}) : rien n'est plié ici, la passe laisse l'arbre
        inchangé, ce qui est correct.
  \item \textbf{Évaluation.} La fonction est d'abord enregistrée (pré-passe), puis
        la boucle \code{pour} appelle \code{factorielle(1)} à
        \code{factorielle(5)}. Chaque appel empile un frame, lie \code{n}, évalue
        le corps, et dépile. Les valeurs retournées sont 1, 2, 6, 24, 120.
  \item \textbf{Sortie.} \code{afficher} imprime chaque résultat :
\end{enumerate}

\begin{lstlisting}[style=gram]
$ ./ilang factorielle.ilang
1
2
6
24
120
\end{lstlisting}

\noindent Le même programme passé à \code{./ilang -s} produit l'assembleur de
l'annexe~\ref{annexe:asm} qui, compilé et exécuté, affiche exactement la même
chose. La boucle complète, depuis le texte source jusqu'au résultat et par deux
back-ends indépendants, est ainsi vérifiée de bout en bout.

\chapter{Discussion}
\label{chap:discussion}

\section{Analyse critique}

\subsection{Forces}

Le projet atteint l'ensemble de ses objectifs et présente plusieurs qualités
notables. Sur le plan de l'\textbf{ingénierie}, la compilation est propre (aucun
avertissement, aucun conflit grammatical) et la gestion mémoire est exempte de
fuite. Sur le plan de la \textbf{conception}, l'architecture en pipeline avec
front-end commun s'est révélée juste : elle a permis d'ajouter un back-end de
génération de code sans toucher à l'interpréteur, validant a posteriori le choix
d'un AST comme représentation pivot. Sur le plan de la \textbf{robustesse}, toutes
les erreurs prévues sont détectées et localisées, et la campagne de cas limites
n'a révélé aucun plantage non maîtrisé.

\subsection{Faiblesses et limitations}

L'honnêteté impose de recenser les limites, dont certaines sont des choix assumés
et d'autres des bornes de fond de l'approche.

\begin{description}[style=nextline,leftmargin=1em]
  \item[Profondeur de récursion bornée.] L'évaluateur récursant en C, la
        profondeur d'appels \ilang{} est limitée par la pile de l'hôte (rupture
        autour de quelques milliers de niveaux, cf. tableau~\ref{tab:rec}). Une
        machine virtuelle à pile explicite lèverait cette limite.
  \item[Performance de l'interprétation.] Le parcours d'arbre est intrinsèquement
        lent (cf. figure~\ref{fig:bench-fib}). C'est le prix de la lisibilité ; un
        bytecode l'améliorerait sensiblement.
  \item[Absence de portée de bloc.] Les variables déclarées dans un bloc
        (\code{si}, boucle) vivent dans le frame de la fonction, pas dans un frame
        de bloc. C'est cohérent avec la spécification (« globale au principal,
        locale aux fonctions ») mais diffère de langages à portée de bloc.
  \item[Back-end assembleur restreint aux entiers.] Le générateur de code refuse
        les flottants et ne reproduit pas la gestion « douce » de la division par
        zéro (un \code{idiv} par zéro lève un SIGFPE). Choix assumé, hors
        périmètre.
  \item[Pas de système de types.] Le langage est dynamiquement typé au sens
        minimal (entier ou flottant), sans vérification statique.
\end{description}

\subsection{Comparaison avec l'existant}

Par rapport à l'approche antérieure (le mini-langage \textsc{Parlot} et son
analyse manuelle en Python), \ilang{} gagne sur deux tableaux : la grammaire est
formalisée et vérifiable, et l'usage d'outils générateurs élimine des classes
entières d'erreurs (oubli de cas, ambiguïté non détectée). En contrepartie, la
courbe d'apprentissage de Flex et Bison est plus raide qu'une descente récursive
écrite à la main. C'est toutefois un investissement qui se rentabilise dès que la
grammaire se complexifie, et qui paie immédiatement sur le plan de la fiabilité,
puisqu'il supprime des familles entières d'erreurs au lieu de les laisser à la
vigilance du programmeur.

\section{Difficultés rencontrées et solutions}

Cette section revient sur les obstacles techniques marquants et la manière dont
ils ont été levés.

\subsection{Le « dangling else » qui n'existait pas}

Le conflit syntaxique du \emph{dangling else}, c'est-à-dire la question de savoir à
quel \code{si} rattacher un \code{sinon} dans une imbrication, est l'exemple
canonique de conflit shift/reduce enseigné dans tous les cours de compilation.
M'attendant à le rencontrer, j'avais d'abord protégé ma grammaire par une
résolution de priorité, en introduisant un pseudo-token de plus faible priorité
que \code{SINON}. Par acquit de conscience, j'ai ensuite retiré cette protection
pour vérifier si le conflit existait réellement, et Bison n'en a signalé
absolument aucun. Mes deux règles, données ci-dessous, se compilent sans la
moindre ambiguïté.

\begin{lstlisting}[style=gram,caption={Les deux règles du \code{si}, sans aucune directive de priorité.}]
si_alors
    : SI '(' expression ')' bloc                { $$ = new_if($3, $5, NULL); }
    | SI '(' expression ')' bloc SINON bloc     { $$ = new_if($3, $5, $7); }
    ;
\end{lstlisting}

L'explication tient entièrement à un détail de ma grammaire que je n'avais pas
mesuré au départ : les corps de conditionnelle sont \emph{obligatoirement} placés
entre accolades, car la règle \code{bloc} l'impose. Or, dans le cours, le conflit
naît de corps réduits à une instruction nue. Ici, comme tout corps se termine par
une accolade fermante, un \code{sinon} est nécessairement précédé du symbole
\code{\}}. Au moment précis où l'analyseur pourrait hésiter, après avoir reconnu
\code{SI ( expression ) bloc}, le symbole de prévision n'est jamais \code{SINON}
mais toujours \code{\}} ou un début d'instruction. Autrement dit, \code{SINON}
n'appartient pas à l'ensemble \textsc{Follow} de la réduction du \code{si} sans
\code{sinon}, et l'ambiguïté n'a tout simplement pas lieu d'être. J'ai donc
supprimé la directive de priorité devenue inutile. Je retiens de cet épisode une
leçon de méthode : il vaut mieux \emph{tester} une hypothèse avec l'outil que
recopier par réflexe une solution apprise, car le contexte précis de ma grammaire
rendait cette solution superflue.

\subsection{L'ordre d'évaluation et la portée}

La gestion des fonctions a révélé deux subtilités que je n'avais pas anticipées et
qui, toutes deux, touchent à la portée des variables.

La première est une question d'\textbf{ordre}. Lorsqu'on appelle une fonction,
il faut impérativement évaluer les arguments dans l'environnement de l'appelant
\emph{avant} d'empiler le cadre de la fonction appelée. Dans ma toute première
version, j'avais inversé ces deux étapes, et une expression aussi banale que
\code{fibonacci(n - 1)} échouait : une fois le nouveau cadre empilé, la variable
\code{n} était cherchée dans ce cadre encore vide au lieu du cadre de l'appelant.
La correction a consisté à séparer nettement les phases, comme le montre l'extrait
suivant.

\begin{lstlisting}[style=c,caption={L'ordre correct : évaluer les arguments avant d'empiler le cadre.}]
/* 1) arguments evalues dans le cadre de l'APPELANT */
for (i = 0; i < node->nargs; i++) argvals[i] = eval(node->args[i]);
/* 2) seulement maintenant on empile le cadre de l'appele */
env_push();
for (i = 0; i < decl->nargs; i++) env_set(decl->args[i]->name, argvals[i]);
\end{lstlisting}

La seconde subtilité, plus profonde, concerne la \textbf{nature de la portée}. Ma
première implémentation de la recherche de variable parcourait toute la pile
d'appel, du sommet jusqu'à la base. Or ce parcours complet réalise en fait une
portée \emph{dynamique} : une fonction pouvait y lire les variables locales de la
fonction qui l'avait appelée, ce qui n'est pas le comportement attendu d'un langage
moderne. Le test reproducteur ci-dessous affichait \code{123} au lieu d'échouer,
puisque \code{g} voyait le \code{secret} local de \code{f}.

\begin{lstlisting}[style=ilang,caption={Cas qui révèle la portée dynamique (corrigé depuis).}]
fonction g() { retourner(secret) }
fonction f() { secret = 123 retourner(g()) }
afficher(f())          // affichait 123 : g voyait une locale de f !
\end{lstlisting}

La spécification réclamant une portée « lexicale », j'ai corrigé la résolution
pour qu'elle ne consulte que le cadre courant, puis en repli l'environnement
global, sans traverser les cadres intermédiaires de la pile d'appel.

\begin{lstlisting}[style=c,caption={Résolution lexicale corrigée (extrait de \code{symtab.c}).}]
int env_get(const char *name, Value *out) {
    if (lookup_in(current_env, name, out)) return 1;          /* cadre courant */
    if (current_env != global_env && lookup_in(global_env, name, out)) return 1;
    return 0;                                                 /* sinon : absente */
}
\end{lstlisting}

Désormais le programme ci-dessus signale correctement que \code{secret} n'est pas
défini dans \code{g}. Je relève au passage une contradiction du cahier des charges
lui-même, qui qualifiait de « lexicale » une portée dont il décrivait pourtant
l'algorithme dynamique. J'ai tranché en faveur du terme employé, qui correspond
aussi au comportement le plus utile.

\subsection{Le piège du pliage de la division par zéro}

Le pliage de constantes consiste à calculer dès la compilation les
sous-expressions dont les deux opérandes sont des littéraux. Le piège est que
\code{10 / 0} est, lui aussi, une expression à deux opérandes littéraux. Si je le
pliais sans précaution, deux issues fâcheuses étaient possibles : ou bien
l'optimiseur lui-même provoquait une division entière par zéro et un
\code{SIGFPE} au moment de la compilation, ou bien il produisait une valeur
arbitraire qui aurait court-circuité le message d'erreur que l'évaluateur doit
émettre à l'exécution, message exigé par le cahier des charges. La solution
consiste à interdire le pliage des seules opérations qui échoueraient, en laissant
intact le sous-arbre concerné pour que l'évaluateur le rencontre normalement à
l'exécution. C'est le rôle du test \code{can\_fold}, invoqué avant tout pliage :

\begin{lstlisting}[style=c,caption={Le pliage n'est tenté que si l'opération est sûre (extrait de \code{optim.c}).}]
if (node->left->type == NODE_NUM && node->right->type == NODE_NUM
    && can_fold(node->op, node->left->num, node->right->num)) {
    Value v = do_fold(node->op, node->left->num, node->right->num);
    /* ... le noeud BINOP est remplace par le litteral resultat ... */
}
/* sinon : on laisse le noeud tel quel, l'evaluateur tranchera a l'execution */
\end{lstlisting}

Ce petit garde-fou illustre une idée importante : une optimisation ne doit jamais
changer le comportement observable d'un programme, pas même ses messages
d'erreur.

\subsection{Détails d'environnement et de robustesse}

Plusieurs petits obstacles, sans gloire mais instructifs, ont jalonné le travail.
Je les mentionne car leur résolution méthodique a directement contribué à la
propreté finale du logiciel.

Le premier est apparu dès la toute première compilation, sous la forme d'un
avertissement déroutant sur la fonction \code{strdup} :

\begin{lstlisting}[style=gram]
implicit declaration of function 'strdup'
assignment to 'char *' from 'int' makes pointer from integer without a cast
\end{lstlisting}

La cause n'est pas une faute de ma part mais une subtilité de la norme :
\code{strdup} ne fait pas partie du C standard, c'est une extension POSIX. Compilé
en \code{-std=c11} strict, l'en-tête masque son prototype, le compilateur suppose
alors un retour de type \code{int}, et sur une machine 64~bits le pointeur se
trouve tronqué. La correction tient en un mot dans le \code{Makefile} :

\begin{lstlisting}[style=gram]
CFLAGS = -Wall -Wextra -std=gnu11 -g    # gnu11 expose les extensions POSIX
\end{lstlisting}

Le deuxième obstacle, encore plus discret, fut une ligne de dépendance du
\code{Makefile} indentée par mégarde avec une espace, là où Make attend soit une
tabulation pour une recette, soit aucune indentation pour une dépendance.
Le troisième fut le tampon interne de Flex, jamais libéré, que Valgrind signalait
comme « encore atteignable » et que l'appel \code{yylex\_destroy()} en fin de
\code{main} a fait disparaître. Le dernier, enfin, fut un débordement silencieux :
le littéral \code{9999999999} s'affichait \code{1410065407} parce que le lexer le
convertissait avec \code{atoi}, limité à 32~bits, alors que mes valeurs sont des
\code{long}. La correction a consisté à remplacer la conversion :

\begin{lstlisting}[style=gram,caption={Avant et après, dans \code{ilang.l}.}]
{CHIFFRE}+   { yylval.num = atoi(yytext); ... }    /* avant : tronque a 32 bits */
{CHIFFRE}+   { yylval.num = strtol(yytext, NULL, 10); ... }  /* apres : 64 bits */
\end{lstlisting}

Aucun de ces points n'est spectaculaire pris isolément, mais c'est l'attention
portée à chacun qui distingue un programme qui « marche sur les exemples » d'un
programme réellement soigné.

\section{L'extension : du parcours d'arbre au code natif}

Le développement du back-end assembleur, bien que hors périmètre, mérite une
discussion car il éclaire la nature même du projet.

\subsection{Modèle et choix}

Le générateur applique la convention cdecl et matérialise, en assembleur, ce que
l'évaluateur réalisait abstraitement : la pile d'appels devient une vraie pile de
cadres \code{\%ebp}, les paramètres et locales des déplacements, les expressions
une suite d'opérations sur \code{\%eax} et la pile. Les structures de contrôle se
traduisent en labels et sauts ; \code{arrêter} et \code{continuer} deviennent
naturellement des \code{jmp} vers les labels de fin et de continuation de boucle.

\subsection{Obstacles spécifiques}

Trois difficultés propres à la cible native ont dû être traitées : l'installation
de la chaîne 32~bits (\code{gcc-multilib}) ; un avertissement de l'éditeur de
liens (\code{DT\_TEXTREL}) dû à l'adressage absolu incompatible avec le PIE,
résolu en liant avec \code{-no-pie} ; et l'alignement de la pile requis par
\code{printf}, assuré par un \code{andl \$-16, \%esp} en tête de \code{main} et
l'arrondi des cadres à 16~octets.

\subsection{Enseignement}

Voir le \emph{même} AST se transformer tantôt en valeurs (interprétation), tantôt
en instructions machine (compilation), rend tangible toute la théorie de la
chaîne de compilation. C'est, de l'avis de l'auteur, la partie la plus formatrice
du projet, et cela tient sans doute précisément au fait qu'elle dépassait la
commande.

\section{Compromis effectués}

En synthèse, les principaux compromis assumés sont : la \emph{lisibilité} au
détriment de la performance, avec un parcours d'arbre plutôt qu'une machine à
bytecode ; la \emph{simplicité} au détriment de l'exhaustivité, sans portée de
bloc, sans typage statique et avec un back-end limité aux entiers ; et la
\emph{fidélité à la spécification} partout où elle était claire. Là où le cahier
des charges se contredisait lui-même, j'ai fait des choix explicites et
documentés. Pour la portée, j'ai retenu le sens lexical conforme au terme employé
plutôt que l'algorithme dynamique décrit. Pour le point-virgule, présent dans la
BNF mais absent des programmes de démonstration, j'ai préféré le rendre optionnel
afin d'honorer les deux sources à la fois plutôt que d'en sacrifier une. Ces choix
sont cohérents avec la finalité pédagogique du projet.

\chapter{Conclusion et perspectives}
\label{chap:conclusion}

\section{Synthèse des contributions}

Ce projet a abouti à un interpréteur complet et fonctionnel pour le mini-langage
impératif \ilang. L'ensemble de la chaîne demandée a été réalisé : analyse
lexicale par Flex, analyse syntaxique LALR(1) par Bison avec construction d'un
arbre syntaxique abstrait, table des symboles gérant la portée par une pile
d'environnements, deux passes d'optimisation, et un évaluateur récursif gérant
variables, expressions, structures de contrôle, fonctions récursives et
entrées-sorties. La gestion des erreurs couvre toutes les catégories du cahier des
charges, chacune localisée à la ligne.

À cette base se sont ajoutées trois extensions, à savoir les nombres flottants,
les instructions de rupture de boucle \code{arrêter} et \code{continuer}, et la
localisation systématique des erreurs. S'y est jointe une exploration assumée
hors du périmètre demandé, un second back-end générant du code assembleur i386,
dont les sorties coïncident exactement avec celles de l'interpréteur. La qualité d'ingénierie visée a été tenue : zéro
avertissement, zéro conflit grammatical, zéro fuite mémoire.

\section{Prolongements possibles}

Plusieurs directions permettraient d'enrichir le projet, et l'architecture
retenue rend la plupart d'entre elles abordables sans refonte. Je les présente en
indiquant à chaque fois ce qu'elles impliqueraient concrètement, car c'est cette
estimation du coût qui distingue une vraie perspective d'un simple v\oe u.

\paragraph{Chaînes de caractères.} Ajouter un type \code{string} serait sans doute
l'extension la plus naturelle après les flottants, car elle suit exactement le
même schéma. Il suffirait d'enrichir l'union étiquetée \code{Value} d'une
troisième variante, et de gérer l'allocation et la libération des chaînes :

\begin{lstlisting}[style=c,caption={Esquisse d'extension de \code{Value} aux chaînes.}]
typedef enum { VAL_INT, VAL_FLOAT, VAL_STR } ValType;
typedef struct Value {
    ValType type;
    long    ival;
    double  fval;
    char   *sval;   /* nouveau : pointeur sur la chaine (a liberer) */
} Value;
\end{lstlisting}

La vraie difficulté ne serait pas le type lui-même mais la gestion de la durée de
vie des chaînes, qui demanderait soit un comptage de références, soit un
ramasse-miettes, pour éviter fuites et libérations prématurées.

\paragraph{Tableaux.} Introduire des tableaux, avec une syntaxe du type
\code{tab = [1, 2, 3]} et un accès indexé \code{tab[i]}, demanderait deux nouveaux
types de n\oe uds dans l'AST, l'un pour le littéral de tableau et l'autre pour
l'indexation, ainsi qu'une nouvelle variante de \code{Value} pointant vers un bloc
alloué. L'évaluateur s'enrichirait de deux cas dans son \code{switch}, ce qui
reste dans la logique du parcours d'arbre existant.

\paragraph{Machine virtuelle à bytecode.} Pour dépasser les deux principales
limites de l'interpréteur, la lenteur du parcours d'arbre et la profondeur de
récursion bornée, la voie classique consiste à compiler l'AST vers un jeu
d'instructions de bas niveau, exécuté par une boucle dédiée disposant de sa propre
pile explicite. On définirait d'abord un jeu d'instructions :

\begin{lstlisting}[style=c,caption={Esquisse d'un jeu d'instructions de bytecode.}]
typedef enum {
    OP_PUSH, OP_LOAD, OP_STORE,      /* constantes et variables */
    OP_ADD, OP_SUB, OP_MUL, OP_DIV,  /* arithmetique sur la pile */
    OP_JMP, OP_JZ,                   /* sauts (controle de flot) */
    OP_CALL, OP_RET, OP_PRINT        /* fonctions et affichage */
} ByteOp;
\end{lstlisting}

L'évaluation reviendrait alors à une boucle lisant ces instructions une à une.
C'est l'approche de Python ou de Lua, et elle supprimerait du même coup la limite
de récursion, puisque la pile d'appels deviendrait une structure de données gérée
explicitement plutôt que la pile du langage hôte.

\paragraph{Back-end assembleur complet, vérification de types, afficheur d'AST.}
Trois autres prolongements me semblent intéressants. Le back-end assembleur
pourrait être étendu aux flottants, au prix d'une gestion du jeu d'instructions
flottantes du processeur, et reproduire la gestion « douce » de la division par
zéro en émettant un test avant chaque division. Une passe d'analyse sémantique
préalable pourrait signaler davantage d'erreurs avant même l'exécution, par
exemple l'usage d'une variable jamais affectée. Enfin, un afficheur d'AST
reconstituant le source avec un parenthésage minimal serait un outil précieux,
aussi bien pour le débogage que pour l'enseignement, car il rendrait visible la
structure que l'analyseur a réellement construite.

\section{Applications potentielles}

Au-delà de l'exercice, les techniques mises en \oe uvre dépassent largement la
construction de langages généralistes. Elles s'appliquent à tout traitement de
texte structuré : langages de configuration, langages dédiés (\emph{DSL}) métier,
moteurs de calcul de formules (tableurs), analyseurs de protocoles, ou encore
préprocesseurs et transpileurs. La compétence centrale, qui consiste à décrire
formellement une syntaxe puis à lui associer une sémantique, est directement
réutilisable dans tous ces contextes, bien au-delà du seul cas des langages de
programmation généralistes.

\section{Recommandations}

Pour qui reprendrait ou referait un tel projet, trois recommandations se dégagent
de l'expérience. Premièrement, \textbf{construire incrémentalement} : un c\oe ur
minimal qui fonctionne de bout en bout vaut mieux qu'un édifice complet jamais
exécuté. Deuxièmement, \textbf{tester tôt et systématiquement}, y compris les cas
limites et la mémoire (Valgrind), plutôt que de se fier aux seuls exemples
heureux. Troisièmement, \textbf{exploiter les outils} plutôt que de les
contourner : laisser Bison signaler les conflits, laisser le compilateur signaler
les avertissements, et traiter chacun comme une information à comprendre.

\section{Apports personnels et conclusion}

Sur le plan des apprentissages, ce projet a transformé des notions jusque-là
théoriques en savoir-faire concret. La distinction entre niveaux régulier et
hors-contexte, la mécanique d'une grammaire LALR(1), le rôle d'un arbre syntaxique
comme représentation pivot, la gestion d'une pile d'appels : autant de concepts
qui prennent un sens entièrement nouveau une fois qu'on les a implémentés et
débogués soi-même, plutôt que seulement lus dans un cours. Le moment le plus marquant aura été de voir le même arbre produire,
selon le back-end choisi, soit un résultat calculé, soit des instructions
machine : cette dualité résume à elle seule ce que le cours cherche à transmettre.

En définitive, \ilang{} n'est pas un langage destiné à être utilisé, mais un
langage destiné à être \emph{compris}. C'est à cette aune qu'il faut juger le
projet : non à la richesse de ses fonctionnalités, mais à la clarté avec laquelle
il rend visible le chemin qui mène d'un texte source à son exécution.


% ===================== Bibliographie =====================
\chapter*{Bibliographie et webographie}
\addcontentsline{toc}{chapter}{Bibliographie et webographie}

\begin{thebibliography}{99}

\bibitem{dragon}
A.~V.~Aho, M.~S.~Lam, R.~Sethi, J.~D.~Ullman,
\emph{Compilers: Principles, Techniques, and Tools} (le « Dragon Book »),
2\textsuperscript{e}~édition, Pearson/Addison-Wesley, 2006.

\bibitem{levine}
J.~Levine,
\emph{flex \& bison: Text Processing Tools},
O'Reilly Media, 2009.

\bibitem{flex}
\emph{Flex : The Fast Lexical Analyzer Generator}, manuel de référence.
\url{https://github.com/westes/flex} (consulté en 2026).

\bibitem{bison}
\emph{GNU Bison Manual}, Free Software Foundation.
\url{https://www.gnu.org/software/bison/manual/} (consulté en 2026).

\bibitem{knuth}
D.~E.~Knuth,
\emph{On the Translation of Languages from Left to Right},
Information and Control, vol.~8, n\textsuperscript{o}~6, p.~607--639, 1965.

\bibitem{chomsky}
N.~Chomsky,
\emph{Three Models for the Description of Language},
IRE Transactions on Information Theory, vol.~2, n\textsuperscript{o}~3, 1956.

\bibitem{lesk}
M.~E.~Lesk, E.~Schmidt,
\emph{Lex : A Lexical Analyzer Generator},
Bell Laboratories, Computing Science Technical Report n\textsuperscript{o}~39, 1975.

\bibitem{johnson}
S.~C.~Johnson,
\emph{Yacc: Yet Another Compiler-Compiler},
Bell Laboratories, Computing Science Technical Report n\textsuperscript{o}~32, 1975.

\bibitem{isoc}
ISO/IEC 9899:2011, \emph{Programming languages : C} (norme C11).

\bibitem{intel}
Intel Corporation,
\emph{Intel 64 and IA-32 Architectures Software Developer's Manual},
volume~2 (Instruction Set Reference).

\bibitem{sysv}
\emph{System V Application Binary Interface : Intel386 Architecture Processor
Supplement} (convention d'appel cdecl).

\bibitem{cours}
P.~Kislin-Duval,
\emph{Interprétation \& Compilation}, support de cours, Licence Informatique,
IED, Université Paris~8.

\end{thebibliography}


% ===================== Annexes =====================
\appendix
\chapter{Code source complet}
\label{annexe:code}

Cette annexe reproduit l'intégralité du code source du projet. Les fichiers
générés automatiquement par Flex et Bison (\code{lex.yy.c}, \code{ilang.tab.c},
\code{ilang.tab.h}) ne sont pas reproduits, étant entièrement déduits de
\code{ilang.l} et \code{ilang.y}.

\section{ast.h}
\lstinputlisting[style=c]{../ast.h}

\section{ast.c}
\lstinputlisting[style=c]{../ast.c}

\section{symtab.h}
\lstinputlisting[style=c]{../symtab.h}

\section{symtab.c}
\lstinputlisting[style=c]{../symtab.c}

\section{ilang.l : analyseur lexical (Flex)}
\lstinputlisting[style=gram]{../ilang.l}

\section{ilang.y : analyseur syntaxique (Bison)}
\lstinputlisting[style=gram]{../ilang.y}

\section{eval.c : évaluateur}
\lstinputlisting[style=c]{../eval.c}

\section{optim.c : optimisations}
\lstinputlisting[style=c]{../optim.c}

\section{codegen.c : back-end assembleur (extension)}
\lstinputlisting[style=c]{../codegen.c}

\section{main.c : point d'entrée}
\lstinputlisting[style=c]{../main.c}

\section{Makefile}
\lstinputlisting{../Makefile}

\section{Programmes de démonstration}

\subsection{fibonacci.ilang}
\lstinputlisting[style=ilang]{../programmes/fibonacci.ilang}

\subsection{pgcd.ilang}
\lstinputlisting[style=ilang]{../programmes/pgcd.ilang}

\subsection{interactif.ilang}
\lstinputlisting[style=ilang]{../programmes/interactif.ilang}

\subsection{flottants.ilang (extension)}
\lstinputlisting[style=ilang]{../programmes/flottants.ilang}

\subsection{boucles.ilang (extension)}
\lstinputlisting[style=ilang]{../programmes/boucles.ilang}

\chapter{Manuel d'installation}
\label{annexe:install}

\section{Prérequis}

Une distribution GNU/Linux (ou WSL sous Windows) disposant de la chaîne C, de
Flex et de Bison. Sur une base Debian/Ubuntu :

\begin{lstlisting}[style=gram]
sudo apt update
sudo apt install build-essential bison flex make
# Optionnel, pour l'extension assembleur 32 bits :
sudo apt install gcc-multilib
\end{lstlisting}

\section{Compilation}

Depuis le répertoire du projet :

\begin{lstlisting}[style=gram]
make          # compile l'interpreteur -> binaire ./ilang
make clean    # supprime les fichiers generes
make test     # compile puis lance tous les programmes de demonstration
make asm      # demonstration du back-end assembleur (extension)
\end{lstlisting}

La compilation ne doit produire aucun avertissement ni conflit.

\chapter{Guide utilisateur}
\label{annexe:guide}

\section{Invocation}

\begin{lstlisting}[style=gram]
./ilang programme.ilang        # execute un fichier source
./ilang                        # mode interactif (Ctrl-D pour executer)
echo 7 | ./ilang interactif.ilang   # alimenter lire() via un tube
./ilang -s programme.ilang     # genere l'assembleur i386 (extension)
./ilang --no-opt programme.ilang    # desactive les optimisations
\end{lstlisting}

\section{Aide-mémoire du langage}

\begin{table}[ht]
\centering
\begin{tabular}{@{}ll@{}}
\toprule
\textbf{Construction} & \textbf{Syntaxe} \\
\midrule
Affectation        & \code{x = 42} \\
Flottant           & \code{r = 2.5} \\
Arithmétique       & \code{+ - * / \%}, \code{-x}, parenthèses \\
Comparaisons       & \code{== != < > <= >=} \\
Condition          & \code{si (c) \{ ... \} sinon \{ ... \}} \\
Boucle tant que    & \code{tant que (c) \{ ... \}} \\
Boucle pour        & \code{pour (i=0; i<n; i=i+1) \{ ... \}} \\
Rupture de boucle  & \code{arrêter} , \code{continuer} \\
Fonction           & \code{fonction f(a,b) \{ ... retourner(e) \}} \\
Affichage / saisie & \code{afficher(e)} , \code{lire()} \\
Commentaires       & \code{// ligne} , \code{/* bloc */} \\
\bottomrule
\end{tabular}
\end{table}

\chapter{Données expérimentales}
\label{annexe:data}

Les mesures ont été obtenues sur AMD Ryzen~5 5500, sous Ubuntu (WSL), meilleur de
trois exécutions. Elles sont reproductibles via les scripts de mesure fournis.

\begin{table}[ht]
\centering
\caption{Données brutes des benchmarks.}
\begin{tabular}{@{}lll@{}}
\toprule
\textbf{Expérience} & \textbf{Paramètre} & \textbf{Mesure} \\
\midrule
Fibonacci (interpréteur) & $n=28$ & 0,43 s \\
Fibonacci (interpréteur) & $n=30$ & 1,13 s \\
Fibonacci (interpréteur) & $n=32$ & 2,97 s \\
Fibonacci (asm i386)     & $n=32$ & 0,01 s \\
Boucle 30~M (avec optim.) & pliage actif   & 4,05 s \\
Boucle 30~M (sans optim.) & \code{--no-opt} & 6,02 s \\
Récursion \code{somme}    & $n=5000$  & ok \\
Récursion \code{somme}    & $n=10000$ & débordement de pile \\
\bottomrule
\end{tabular}
\end{table}

\chapter{Assembleur i386 généré (exemples)}
\label{annexe:asm}

Cette annexe reproduit le code assembleur réellement produit par le back-end
(\code{./ilang -s}) pour deux programmes. Ce code est compilable tel quel avec
\code{gcc -m32 -no-pie} et produit, à l'exécution, exactement les mêmes résultats
que l'interpréteur (respectivement 6 et 120).

\section{PGCD (boucle \code{tant que})}
Source : voir \code{pgcd.ilang} en annexe~\ref{annexe:code}. Le code généré
illustre la traduction d'une boucle en labels et sauts conditionnels.
\lstinputlisting[style=asm]{generated/pgcd.s}

\section{Factorielle (fonction récursive)}
Pour \code{factorielle(5)}, le code généré illustre la récursion : la fonction
\code{il\_factorielle} s'appelle elle-même, chaque appel disposant de son propre
cadre de pile.
\lstinputlisting[style=asm]{generated/factorielle.s}

\chapter{Contenu de l'archive}
\label{annexe:archive}

Le code source est livré dans l'archive
\code{IED-L3-IC-CHAPITRE\_10-ILIAS\_AIT\_BENAISSA-82508880-sources.zip}, dont
l'arborescence est la suivante :

\begin{lstlisting}[style=gram]
ilang/
|-- ilang.l            analyseur lexical (Flex)
|-- ilang.y            grammaire et construction de l'AST (Bison)
|-- ast.h  ast.c       definition et manipulation de l'AST, type Value
|-- symtab.h  symtab.c table des symboles (pile de cadres + hachage)
|-- eval.c             evaluateur recursif
|-- optim.c            optimisations (pliage de constantes, code mort)
|-- codegen.c          back-end assembleur i386 (extension hors-perimetre)
|-- main.c             point d'entree (fichier, interactif, -s, --no-opt)
|-- Makefile           compilation automatisee
|-- README.md          instructions de compilation et d'utilisation
|-- JOURNAL.md         journal de bord (difficultes, choix de conception)
|-- programmes/        programmes de demonstration
|   |-- fibonacci.ilang
|   |-- pgcd.ilang
|   |-- interactif.ilang
|   |-- flottants.ilang     (extension : flottants)
|   |-- boucles.ilang       (extension : arreter / continuer)
|-- tests/             cas de test
    |-- err_variable.ilang   err_divzero.ilang   err_fonction.ilang
    |-- err_arguments.ilang  err_lexical.ilang   err_syntaxe.ilang
    |-- optim.ilang          run_errors.sh
\end{lstlisting}

\paragraph{Remarque sur les fichiers générés.} L'archive ne contient
volontairement \emph{pas} les fichiers produits automatiquement lors de la
compilation, à savoir \code{lex.yy.c} (engendré par Flex à partir de
\code{ilang.l}), ainsi que \code{ilang.tab.c} et \code{ilang.tab.h} (engendrés par
Bison à partir de \code{ilang.y}). Ces fichiers sont entièrement reconstruits par
la commande \code{make} et n'ont donc pas à figurer dans les sources. De même, les
fichiers objets \code{.o} et le binaire \code{ilang} ne sont pas livrés, puisqu'ils
résultent de la compilation.

\paragraph{Récapitulatif des sorties.} Le tableau~\ref{tab:recap} rassemble, pour
mémoire, les sorties attendues des programmes de démonstration, déjà détaillées au
chapitre~\ref{chap:validation}.

\begin{table}[ht]
\centering
\caption{Récapitulatif des sorties des programmes de démonstration.}
\label{tab:recap}
\small
\begin{tabular}{@{}lll@{}}
\toprule
\textbf{Programme} & \textbf{Entrée} & \textbf{Sortie} \\
\midrule
\code{fibonacci.ilang}  & aucune & \code{0 1 1 2 3 5 8 13 21 34 55 89 144 233 377} \\
\code{pgcd.ilang}       & aucune & \code{6} \\
\code{interactif.ilang} & \code{7}  & \code{999} puis \code{49} \\
\code{interactif.ilang} & \code{-5} & \code{999} puis \code{5} \\
\code{flottants.ilang}  & aucune & \code{6.28 / 3.5 / 2 / 2.5 / 19.6349 / 1} \\
\code{boucles.ilang}    & aucune & \code{0 1 2 3 4 / 1 3 5 7 9 / 10 9 8} \\
\bottomrule
\end{tabular}
\end{table}

\chapter{Glossaire}
\label{annexe:glossaire}

\begin{description}[style=nextline,leftmargin=1em]
  \item[AST] \emph{Abstract Syntax Tree}, arbre syntaxique abstrait :
        représentation arborescente du programme, débarrassée des détails de
        syntaxe concrète.
  \item[BNF] \emph{Backus-Naur Form}, notation pour décrire formellement une
        grammaire hors-contexte.
  \item[cdecl] convention d'appel de fonction sur architecture x86 (arguments
        empilés de droite à gauche, nettoyage par l'appelant).
  \item[LALR(1)] \emph{Look-Ahead LR}, classe de grammaires analysables de manière
        ascendante avec un symbole de prévision ; celle qu'accepte Bison.
  \item[Lexème / token] unité lexicale élémentaire reconnue par l'analyseur
        lexical (nombre, identifiant, mot-clé, opérateur).
  \item[PIE] \emph{Position-Independent Executable}, format d'exécutable par
        défaut sur les systèmes modernes, incompatible avec l'adressage absolu
        utilisé par le back-end (d'où l'option \code{-no-pie}).
  \item[Pliage de constantes] optimisation calculant à la compilation les
        sous-expressions entièrement constantes.
  \item[Portée lexicale] règle de visibilité des variables fondée sur la structure
        du programme (et non sur la chaîne dynamique des appels).
  \item[Shift/reduce] type de conflit qu'un analyseur LR peut rencontrer lorsqu'il
        hésite entre décaler un symbole et réduire une règle.
\end{description}


\end{document}
